\documentclass[10pt,letterpaper]{article}
\usepackage[margin=0.85in]{geometry}
\usepackage[T1]{fontenc}
\usepackage{lmodern,microtype,amsmath,amssymb,amsthm,mathtools}
\usepackage{booktabs,longtable,array,enumitem,xcolor,xurl,hyperref,fancyvrb}
\hypersetup{colorlinks=true,linkcolor=blue!45!black,urlcolor=blue!45!black,citecolor=blue!45!black,
pdftitle={A Start-only marker and compositional reductions for Diophantine consistency of ZFC},
pdfauthor={Research prepared for Vladimir Reshetnikov and ProveIt with ChatGPT assistance},
pdfsubject={Explicit Diophantine consistency equations, ordinary proofs, and reproducible coefficient certificates}}
\setlength{\parindent}{1.1em}
\setlength{\parskip}{0.25em}
\setlength{\emergencystretch}{2em}
\setlist{nosep,leftmargin=*}
\newcommand{\Z}{\mathbb Z}
\newcommand{\N}{\mathbb N}
\newcommand{\Pos}{\mathbb Z_{>0}}
\newcommand{\Con}{\operatorname{Con}}
\newcommand{\pc}{\operatorname{pc}}
\newcommand{\lcm}{\operatorname{lcm}}
\newcommand{\sgn}{\operatorname{sgn}}
\newtheorem{theorem}{Theorem}[section]
\newtheorem{lemma}[theorem]{Lemma}
\newtheorem{proposition}[theorem]{Proposition}
\theoremstyle{remark}\newtheorem{remark}[theorem]{Remark}
\title{A Start-only marker and compositional reductions\\for Diophantine consistency of ZFC}
\author{Research prepared for Vladimir Reshetnikov and ProveIt\\\small Mathematical development and checking assisted by ChatGPT}
\date{25 October 2026}
\begin{document}
\maketitle
\begin{abstract}
We give three explicit improvements to the constructions in ProveIt at the
fixed comparison commit \texttt{fb1bed5eb7196cf35abfaa746c18c6ee7eae6cfb}.
First, a physical flag can track the predecessor of the distinguished Start
tile instead of tracking every occupied cell. A nine-position alignment
argument removes 942 terms from a constant-defining kernel and reduces a
native cell from 6007 to 4127 radix digits. The four fixed integer
coefficients then occupy 97,046 minimal unsigned bytes rather than 141,226,
with the same 15 positive variables, 66 arithmetic operations and degree 157.
Second, combining an explicit Pell trace with an inverse Lucas correction
gives a ten-positive-variable equation using 85 operations at exact degree
232. A square-transport variant also settles a previously stated local open
step and reduces its carrier degree from 57 to 21.
Third, a linear binomial selector at a self-power horizon, followed by exact
integer projections, gives a quartic in 233 unrestricted integer variables
with 169 quadratic residuals. Its balanced form has coefficient height 5632.
All equations and fixed constants are supplied explicitly. The new proofs
compose with identified ordinary source calibrations to prove, in ZFC, that
unsolvability in each stated domain is equivalent to $\Con(\mathrm{ZFC})$.
Finite artifact checks and inherited unbounded proofs are distinguished;
no global minimality, worldwide priority, or new end-to-end formal-kernel
certificate is asserted.
\end{abstract}
\setcounter{tocdepth}{1}
\tableofcontents
\clearpage

\section{Results and the meaning of ``simplest''}

The requested objective has several independent size measures. A polynomial
with fewer variables can require much higher degree or very large fixed
coefficients; a small evaluation circuit need not have a short expanded
monomial list. We therefore give explicit improvements in three relevant
directions rather than assert a global minimum for an unspecified simplicity measure.
The domain is part of every theorem.

\begin{center}\small
\begin{tabular}{>{\raggedright\arraybackslash}p{3.3cm}rrrr>{\raggedright\arraybackslash}p{3.4cm}}
\toprule
Construction & Variables & Degree & Gates & Height & Main improvement\\
\midrule
Start-only compact form & $15$ positive &157&66& large &97,046 coefficient bytes\\
Trace/inverse-Lucas form &$10$ positive&232&85& large &87/232 and 85/304 combined\\
Direct balanced quartic &$233$ integer&4&---&5632&169 explicit residuals\\
Direct sparse quartic &$233$ integer&4&---&8880&10,182 collected monomials\\
Sparse companion &$234$ integer&4&---&8880&5,018 collected monomials\\
\bottomrule
\end{tabular}
\end{center}

Here a gate is one binary addition, subtraction or multiplication, including
a multiplication by a fixed scalar. Repeated squaring is charged. Fixed
integer leaves are specified separately by a finite recipe and actual bytes.
Height denotes the largest absolute coefficient of the fully collected
polynomial, not the largest coefficient of a residual. The direct balanced
quartic has 18,482 collected monomials; its 169 residuals have 1822 terms
in total and maximum residual coefficient 40.

The maintained compact equation at the comparison commit used a cell of
6007 radix digits and 141,226 coefficient bytes \cite{current,native}.
The direct 234-variable construction of package 159 had 171 residuals,
19,050 monomials and height 4480 \cite{radical}. Our 233-variable forms
reduce dimension and residual count, with an explicit height tradeoff.
Other repository branches already have fewer than 233 integer variables;
our direct result concerns the small-coefficient C222 branch and does not
supersede those other tradeoffs.

For a fixed ordinary proof predicate $\mathrm{Prf}_{\mathrm{ZFC}}$, put
\[
 \Con(\mathrm{ZFC})\equiv
 \neg\exists p\in\N\;\mathrm{Prf}_{\mathrm{ZFC}}(p,\bot).
\]
Each calibration below proves existence of a polynomial zero equivalent to
existence of a finite syntactic refutation. Negation gives the requested
unsolvability equivalence; it does not prove either side separately.

\section{The compact equation: all 15 unknowns and 66 gates}
\label{sec:compact}

Let the only unknowns be
\[
 J,F,\alpha,z,f,h,i,T,s,w,\tau,\eta,\zeta,y,\gamma\in\Pos.
\]
The following are polynomial abbreviations, not additional unknowns:
\begin{align*}
 q_-&=\beta J,&q&=q_-+1,& C&=q-F-\alpha,\\
 X&=wq,&Y&=sq^3,&E&=XY,\\
 k&=\eta+\zeta,&\ell&=EkY,&c&=kY+\eta,\\
 a&=E+Y,&H&=4a+3,&\Delta&=a^2+H,\\
 D&=X+ac+\gamma H,&&&\\[-2pt]
 R&=[q_-(q-F)+\alpha](q^2-1)+(M_C+qM_F)J,&&&\\
 V&=f^2(cT-R)-c,& S&=i\Delta c^2.&&
\end{align*}
Define six factors
\begin{align*}
 N_0&=\tau^2-\ell(\ell+k),&
 N_1&=D^2-\Delta c^2,\\
 N_a&=S^2(V^2-y^2)+y^2,&
 N_k&=k-hE-R,\\
 N_t&=(K+w)C+(q-F)-zq_-,&
 N_s&=\Delta f^2-S^2.
\end{align*}
The preferred compact equation is
\begin{equation}\label{eq:p66}
 \boxed{P_{66}=N_0N_1N_aN_kN_tN_s-\Delta=0.}
\end{equation}
Its literal graph uses 38 multiplications, 15 additions and 13 subtractions.
The new work changes the complete fixed tuple $(\beta,K,M_C,M_F)$, coherently,
while retaining this graph and its positive-domain arithmetic theorem.

\subsection{A finite definition of the four coefficients}
Put
\[
 \rho=2^{47},\quad M=2^{44},\quad L=4127,\quad
 B=\rho^L=2^{193969},\quad\beta=B-1,
 \qquad G_m(x)=\sum_{j=0}^{m-1}x^j.
\]
Appendix~\ref{app:language} defines the entire fixed list of 943 tiles
$t_j=(w_j,n_j,p_j,s_j)$, including the full 436-character numeral and
decoder. No externally chosen machine index is left unspecified. Define
\begin{align*}
 \Psi(x,t)&=x^{992}(x^{10}+t^{17})G_8(tx)
       +t^8x^{943}(x^{30}+t^{17})G_9(tx)
       +t^{34}(2-t)x^{1031},\\
 \Phi(x)&=\sum_{j=0}^{942}
 [1+2^{10}(1+w_j+2^8n_j+2^{17}p_j+2^{25}s_j)]x^j
       +2^{10}\Psi(x,2),\\
 \Lambda&=\rho^{2042}+\rho^{3074}.
\end{align*}
The four fixed integers are
\begin{align}
 K&=\rho^{20}+\rho^{30}+\rho^{1031}+\rho^{4137}
             +\Lambda\Phi(1/\rho),\label{eq:K}\\
 M_C&=\beta+\rho-(\rho+3)\rho^{990}-G_{943}(\rho)-\Psi(\rho,1),\\
 M_F&=B+3+(M-2)\Lambda+\rho^{973}G_9(\rho)
       +\rho^{1012}G_8(\rho)+\rho^{1031}
       +2^{29}\rho^{4124}-\rho^{4106}.\label{eq:MF}
\end{align}
The reciprocal in \eqref{eq:K} is only finite notation: all its exponents
after reflection are nonnegative. The first reflected nonzero term is at
$2042-1009=1033$. Thus all four leaves are ordinary explicit integers.
$M_F$ is already shifted by $\beta$ once; it must not be shifted again.

\begin{theorem}\label{thm:compact}
For the four integers just specified, ordinary mathematics formalizable in
ZFC proves
\[
 \bigl(\forall\mathbf x\in\Pos^{15},\ P_{66}(\mathbf x)\ne0\bigr)
 \ \longleftrightarrow\ \Con(\mathrm{ZFC}).
\]
The exact degree is 157 and the complete evaluator has 66 live gates.
The four minimal unsigned big-endian coefficient strings total 97,046 bytes.
\end{theorem}

The proof of the new native interface follows in the next section. The
unchanged ordinary P66 arithmetic theorem converts that interface to a
positive zero. Its normalized Pell and signed-recovery arguments are
retained explicitly in the source packet. In particular it retains the
possibility of two negative unit factors until the rank argument resolves
their signs; it does not assume that all factors are positive at the outset.

For the degree statement, the six factors have exact degrees
$22,18,62,7,2,46$. The apparent top cancellation in $N_1$ is accounted for
by writing $Z_0=X+\gamma H$ and
\[
 N_1=2acZ_0+Z_0^2-Hc^2.
\]
Writing $q_0=\beta J$ and $k_0=\eta+\zeta$, the complete top homogeneous
part of the product is
\[
 -8q_0^{91}h\gamma i^4k_0^{13}w^{13}s^{26}T^2f^4
       [w(q_0-F-\alpha)-zq_0].
\]
It is nonzero at the actual $\beta>0$, so degree 157 is attained.

\section{The new idea: the flag tracks Start}
\label{sec:native}

The repository's active \texttt{SIZE\_REDUCTION.md} already suggests trying
a Start-only marker. The contribution here is a complete working layout,
the semantic exclusion of its last false alignment, and a new authenticated
controller for the actual shorter coefficients.

In the immediate predecessor, a block of 943 unit coefficients copies the
number of selected tiles into a physical flag test. Our new singleton term
$\rho^{1031}$ copies only selector 0, which denotes Start. All other cells
can have flag zero. This is sufficient to align the temporal shift because
Start also has eight consecutive prescribed label bits.

The physical support is
\[
 \begin{split}
 \mathcal E={}&\{0,\ldots,951\}\cup\{973,\ldots,981\}
 \cup\{992,\ldots,999\}\cup\{1002,\ldots,1009\}\\
 &\cup\{990,991,1031\}.
 \end{split}
\]
It has 980 positions. Apart from the low dummy at 990, whose digit can be
$0,1,2,3$, physical digits are Boolean. The flag at 1031 has coefficient
zero in both arithmetic clauses. It must remain in the physical support
even though it is absent from the nonzero weight support.

Let $c_e$ denote the coefficients of $\Phi$. They consist of the 943
selector weights in the preceding section and the 34 powers of two for
copied labels. Every $c_e<M$, there are 977 nonzero $c_e$, and the label
weights sum to $M-1024$. The two reflected clause centers are 2042 and 3074.
Their separation, 1032, exceeds the largest physical position, 1031.

\subsection{Clauses are decoded before alignment}

The unchanged local clause lemma uses a selector count $n$, a copied
34-bit word $W$, and complemented tile words $u_j$. Its raw value is
\[
 n+2^{10}\left(2^{34}n+W-\sum_{j\text{ selected}}u_j\right).
\]
With incoming carry between 0 and 29, avoiding the mask $M-2$ forces either
no selector and all-zero labels, or exactly one selector and its correct
complemented labels. The flag is not used by this argument.

For the full cyclic convolution, the two reflected copies cannot both use
the same weight at one output. The extra 0/2 at the low dummy is included
separately in the exact convolution envelope. Its maximum, together with
the spatial contribution and a temporal contribution in any fractional
binary phase, satisfies
\[
 H_{\max}+3+\frac{3\rho}{4}+29<30\rho.
\]
Thus the monotone cyclic carry map preserves the finite box $0,\ldots,29$.
Iteration from zero stabilizes; cyclic telescoping identifies its digit
word with the required transport residue. The inherited strict outer
representative excludes the zero/all-ones ambiguity. Since the clause
centers differ by a missing support difference, at least one clause is
clean in each cell. The local lemma therefore decodes every cell before
alignment is used.

\subsection{The decisive nine-position lemma}

The forced Start tile has complemented south label $510$, whose bits
1 through 8 are one. Its eight copy tests at positions
$974,\ldots,981$, and the singleton Start test at 1031, all require odd
temporal contributions. Below the first reflected coefficient, the
non-temporal digits are bounded by 9, so
\[
 9+3\rho/4+29<\rho
\]
resets every carry relevant to those tests. Fractional binary phases
1 through 45 have only even temporal digits. Phase 46 can produce only
one odd low-dummy spill, which cannot meet all nine targets. The phase
is therefore zero.

Let $u$ be the preimage of position 974 under the remaining radix-digit
shift. The eight consecutive preimages must be physical and must include
at most one selector in each predecessor cell. A direct interval argument
gives exactly
\[
 u\in\{942,943,944,973,974,990,991,992,1002\}.
\]
The ninth preimage is $u+57$. Physical support leaves only $u=942$ and
$u=974$. In the first case the bit at selector 942 must be one, so the
decoded tile is exactly
\[
 t_{942}=(4,437,2,437).
\]
Its complemented south label is $511-437=74$, whose low bit is zero.
But the temporal target requires the physical bit at 943 to be one.
This contradiction eliminates $u=942$. Thus $u=974$ and the whole
radix-digit shift is zero. This final use of one authenticated tile is
what permits moving the flag from the intermediate position 1042 to 1031.

\subsection{Resetting carries and recovering a computation}

After alignment, the 17 ordinary copy tests impose all temporal and
spatial edge equalities. The new flag test imposes exactly
\[
 \mathrm{flag}_{\text{temporal predecessor of }i}
     =\mathrm{StartSelector}_i.
\]
Temporal rotation is a permutation, so these equations uniquely determine
the flags. Occupation follows separately: empty cells have zero colors,
genuine tiles have positive complemented colors, and spatial one-cell
rotation propagates occupation from Start to every cell.

Put $A_{\rm adj}=\sum_{e,e+1\in\mathcal E}c_{e+1}x_e$. The complete
clause predecessor values are
\[
 A_{\rm adj}\quad\text{at 2041},\qquad
 A_{\rm adj}+c_0x_{1031}\quad\text{at 3073}.
\]
The second summand is essential: it comes from the unique difference
$1031-0$. At most one selector and the label-weight bound give
$A_{\rm adj}\le2M-1024$. Since $c_0<M$, adding incoming carry 29 cannot
overflow $\rho$. Both clauses reset. The maximum raw product field is
$3074+1031=4105$; its value is only $c_0x_{1031}<M$, and it also resets.
All fields 4106 through 4126 are raw zero. They contain the 875 padding
tests and a terminal zero field.

Start's north and south colors differ, so the temporal stride is nonzero.
Inverse complementation therefore recovers the unchanged calibrated
marked blank computation and finite refutation. This proves soundness of
the new layout without any controller-prime assumption.

\subsection{Completeness uses a fresh controller}

A finite refutation supplies independently padded one-Start rectangles
by the retained helical theorem. Encode their selectors and complemented
labels on the new support; set each flag according to the displayed
Start equation. The 990-dummy bits remain freely selectable. Every
simultaneous subset preserves the masks and carry resets. In particular
the terminal zero field gives $0<F<q/\rho$, while the support bound gives
$0<C<q/\rho^{L-1032}$. Consequently $2C+F<q/3$ and
$\alpha=q-F-C>0$.

For a safe dummy bit at cell $j$, the ordinary change of the actual rank is
\[
 \delta R=-\Gamma B^j,\qquad
 \Gamma=\rho^{990}(q^2-1)[1+q(K+B^{\mathrm{width}})].
\]
It is an ordinary equality, not only a congruence: every intermediate
subset obeys the same strict representative bounds. The new actual
$K$ is used in every controller calculation. The supplied finite
certificate proves primality, all required exact orders, a valuation-one
lift and nonzero controller residues. The principal residues are
\begin{center}
\begin{tabular}{rr}
\toprule
Prime $p$ & $1+BK+B^2\pmod p$\\\midrule
47&9\\4127&3487\\474491&153928\\\bottomrule
\end{tabular}
\end{center}
Here $474491>193969$, and $474490$ is a multiple of $47449$, the least
common multiple of the two fixed orders. The certificate supplies the
exact nontrivial lift modulo $474491^2$.

The inherited prescribed-dimension padding contract permits fresh width
and height to be sufficiently large powers of this auxiliary prime;
mere abstract cofinality would not be enough here.
The retained elementary orbit/CRT argument then makes
$\Gamma$ a unit modulo the whole rank modulus and expresses every desired
rank residue as a Boolean subset sum of safe dummy increments. This
restores the authentic power rank, not a surrogate index. Explicitly, with
$d=47L$ and $M_F^{\rm native}=M_F-\beta$, the masks obey
\[
 0<M_C,M_F^{\rm native}<\beta,\qquad
 \pc(M_C)=d-980,\qquad\pc(M_F^{\rm native})=980,
\]
\[
 M_C\equiv2\pmod4,\qquad M_F^{\rm native}\equiv4\pmod8.
\]
For a canonical tableau with $N$ cells and $t=dN$, the actual rank satisfies
$\pc(R)=3t+2$ and $R\equiv3\pmod4$.
Together these supply the entire positive native interface.
The ordinary P66 theorem gives all
remaining positive witnesses, completing Theorem~\ref{thm:compact}.

\paragraph{An intermediate with a simpler alignment proof.}
The archive also retains the independently checked flag-1042, cell-4159
construction. It needs only physical support and one-selector information
for the ninth target, without examining tile 942. Its four coefficients
occupy 97,798 bytes. The preferred flag-1031 refinement saves 752 additional
bytes. Both are complete coherent tuples; their coefficients must not be mixed.

\section{Ten positive variables: 85 gates and degree 232}
\label{sec:85}

This result composes the trace construction of package 163 with the inverse
Lucas correction of package 164 \cite{trace,inverse}. Their full four-port
coefficient bindings are byte-identical. The archive includes those actual
four integers, the complete printed recipe in Appendix~\ref{app:85constants}, and a standalone
standard-library generator that reproduces every raw and compressed hash.
Here $\beta=2^{590069}-1$. These are the retained radix-67 coefficients,
not the new tuple of Section~\ref{sec:native}.

The only unknowns are
\[
 J,F,\alpha,i,T,w,s,n,\eta,\zeta\in\Pos.
\]
Define
\begin{align*}
 h&=\beta J,&q&=h+1,& C&=q-F-\alpha,\\
 X&=qw,&Y&=qs,&k&=\eta+\zeta,\\
 c&=kY+\eta,&A&=Y(X+1)+1,&\Delta&=A^2-1,\\
 H&=4A-5,&v&=XY^2,&\ell&=vk=(XY)(kY),\\
 E&=1+\ell(\ell+k),&&&\\
 R&=[h(q-F)+\alpha](q^2-1)+(M_C+qM_F)J,&&&\\
 B_t&=(K_0+w)C-F,&U&=ic^2,&a&=\Delta U^2E,\\
 F_0&=1+a,&G&=EF_0,&f&=a+F_0,\\
 Q&=4\Delta aF_0,&b&=cT-R,&V&=fb-c,\\
 L_t&=bV,&K_t&=X(3c-2X)+2,&&\\
 e_6&=6(k-R-1)-\ell(k^2-4),&D_6&=6q^6,&M_t&=hD_6H,\\
 Z&=hD_6[H(n-L_tB_t)+K_t]
       -H[D_6B_t+h(e_6-6q^4s)].&&&
\end{align*}
The polynomial is
\begin{equation}\label{eq:p85}
 \boxed{P_{85}=Q(Z^2-M_t^2GV^2)+(M_t^2G-Z^2)=0.}
\end{equation}
The literal graph does not compute $v$ separately; $\ell=(XY)(kY)$ reuses
already live registers. Fixed powers are repeated multiplications.

\begin{theorem}
With the included calibrated four-port tuple, ZFC proves absence of
positive integer zeros of \eqref{eq:p85} equivalent to $\Con(\mathrm{ZFC})$.
Its 85 live gates comprise 51 multiplications, 18 additions and 16
subtractions. Its exact degree is 232.
\end{theorem}

\subsection{Why these two reductions can be combined}

For the proof only, write
\[
 e^*=e_6/6,\quad\rho_*=e^*/q^6,\quad\sigma=s/q^2,
 \quad t=B_t/h,\quad\kappa=K_t/H,\quad z=Z/M_t,
 \quad W=1+hL_t.
\]
The exact rational identity is
\begin{equation}\label{eq:carrier}
 z=n+\kappa+\sigma-Wt-\rho_*.
\end{equation}
No quotient in this proof notation is an operation in the polynomial.

The root extraction and signed folding argument of package 163 does not
use the cubic correction. At any positive zero it gives
\[
 E=\tau^2,\quad F_0=f_0^2,\quad G=g^2,\quad
 z=\varepsilon gy,\quad QV^2-(Q-1)y^2=1,
 \qquad\varepsilon\in\{-1,1\}.
\]
It uses the exact identities
\[
 f^2=1+4aF_0,\quad Q=4\Delta^2U^2G,
 \quad f-1=2\Delta i^2c^4E,
\]
prime valuations, and the rational adjacent-Pell descent. Put $I=2Ug$.
Then $f^2-\Delta I^2=1$, $Q=(\Delta I)^2$, and $c^2\mid I$.
Signed folding yields $c=\psi_A(p)$ for odd positive $p$, while the first
Pell norm gives $k=2\psi_{1+2v}(m)$, $m\ge2$ and $c>8XY^3\ge8q^4$.

The main and auxiliary index-three exclusions are unchanged. In particular
an auxiliary index three would force one of
\[
 f(b-4\Delta f)=c-(4\Delta+3),\qquad
 f(b+4\Delta f)=c+(4\Delta+3).
\]
Both right-hand sides are strictly between zero and $f$, a contradiction.
The resulting signed estimates are
\[
 b\ne0,\quad\sgn V=\sgn b,\quad L_t>0,
 \qquad gy<L_t/2,\quad |b|/q^6<L_t/2.
\]
They are proved before the signs of $B_t,C,R,V$ are assumed.

The inverse correction now makes the sign step particularly short:
\[
 e^*=k-R-1-\frac{vk(k^2-4)}6\le k-R-1<b,
\]
because $k\ge2$ and $cT-k+1>0$. If $B_t\le-1$, then $Wt<-L_t$;
\eqref{eq:carrier} would give
\[
 n< L_t/2-L_t+|b|/q^6<0.
\]
Thus $B_t\ge0$, then $C>0$. Only now is the mask bound
$3q<R<q^4-2$ licensed. It gives $c>2R$, hence $b,V>0$.
Since the explicit trace satisfies $f\equiv1\pmod c$, both residues
$V\equiv-c\pmod f$ and $V\equiv-R\pmod c$ survive the trace replacement.
The retained rank theorem gives $p=R$.

\subsection{The denominator of the inverse correction}

Write the exact Lucas expansion
\[
 k=2m+8vC_m+32v^2Z_m,\qquad C_m=\binom{m+1}{3},
 \quad A_m=C_m+4vZ_m,
\]
where $Z_m$ is a nonnegative integer. For $d=2m-R-1$ and
$S_m=k^2+2mk+4m^2-4$, direct algebra gives
\begin{equation}\label{eq:inverseidentity}
 3(e^*-d)=4v^2(24Z_m-A_mS_m).
\end{equation}
Since $k$ is even, $k(k^2-4)/6$ is integral, hence $e^*\in\Z$.
Clearing the integral carrier and using $\gcd(q,hH)=1$ gives
$q^6\mid e^*-q^4s$. Since $v=q^3ws^2$,
the valuation of the right side of \eqref{eq:inverseidentity} at every
prime dividing $q$ gives $q^4\mid e^*-d$, including the prime 3.
Monotonicity of the first Pell base gives $m<R$, so
$|d|<R<q^4$ and $d=0$.

It would be incorrect to infer $q^6\mid e^*$ at this stage. The carrier
first gives $H\mid6hK_t$. The inherited strict whole-product projection,
using the authenticated index, forces $X=2^R$. Since $q\mid X$, $q$ is
a power of two. Only then can the factor 3 be canceled in
\eqref{eq:inverseidentity}. This yields $q^6\mid e^*$, $q^2\mid s$ and
$h\mid B_t$. These are precisely the fixed-minus parity, half-binomial
rounding and native reconstruction interfaces in the retained sources.

\subsection{Completeness is constructive and positive}

At an old package-164 positive zero, put
$d_0=2\Delta i^2c^3E$, so $f=1+cd_0$, and define
\[
 T_{\rm new}=fT_{\rm old}-d_0R,\qquad
 n_{\rm new}=n_{\rm old}+(f-1)L_{\rm old}B_t.
\]
Then $b_{\rm new}=fb_{\rm old}$, $V_{\rm new}=V_{\rm old}$ and
$n_{\rm new}-L_{\rm new}B_t=n_{\rm old}-L_{\rm old}B_t$.
The entire polynomial value is preserved by this polynomial substitution.
The old proof gives $0<R<c/2$, $B_t\ge0$ and $b_{\rm old},V_{\rm old}>0$,
so both new inputs are strictly positive. This supplies completeness.
Soundness was proved separately; an integral inverse to this map is not
assumed.

For exact degree, put $q_0=\beta J$, $k_0=\eta+\zeta$, and
$B_0=k_0sT-q_0^2(q_0-F)$. The full leading homogeneous part is
\[
 [P_{85}]_{232}=-9\cdot2^{10}q_0^{100}w^{26}s^{58}i^{10}k_0^{32}B_0^2.
\]
The unique dominant final term is $-QM_t^2GV^2$, of degree
$80+22+50+80=232$. The displayed leader has 202 collected terms and a
nonzero coefficient $-9216\beta^{100}$ at
$J^{100}w^{26}s^{60}i^{10}\eta^{34}T^2$.

\subsection{A resolved local open step: square transport}

Replace $L_t=bV$ by $L_t=b^2$. The index exclusion already gives
$|b|>8f^3-1/2$, $g<f/2$, $|V|<(3/2)f|b|$, and
$y<(3/2)|V|$; hence $gy<b^2/2$.
The inverse correction still has $e^*<b$, so the signed bootstrap works
without the adverse negative-rank cubic estimates that obstructed this
variant in package 163. Completeness uses its canonical positive auxiliary
pair and the retained inverse-correction bound
$g>\kappa+\sigma+|\rho_*|$ to choose positive $n$ directly.
The variant retains 85 gates and degree 232, but the carrier numerator
$Z$ has exact degree 21 instead of 57. It is included as a proved
structural alternative, not counted as a separate headline Pareto gain.

\section{A direct quartic in 233 unrestricted integers}
\label{sec:direct}

Let $r_0,\ldots,r_{168}\in\Z[x_0,\ldots,x_{232}]$ be the literal quadratic
residuals printed in Appendix~\ref{app:residuals} and supplied as sparse
coefficient arrays. Define
\begin{equation}\label{eq:direct}
 \boxed{P_{233}(x_0,\ldots,x_{232})=
              \sum_{j=0}^{168}r_j(x_0,\ldots,x_{232})^2=0.}
\end{equation}
Every coordinate is unrestricted. In particular no positivity condition
is being silently placed outside this equation.

\begin{theorem}
The integer polynomial \eqref{eq:direct} has exact degree four, 233 live
variables, 169 quadratic residuals, 18,482 collected monomials and height
5632. ZFC proves
\[
 \bigl(\forall\mathbf x\in\Z^{233},\ P_{233}(\mathbf x)\ne0\bigr)
       \longleftrightarrow\Con(\mathrm{ZFC}).
\]
\end{theorem}

Its arithmetic source is the direct C222 construction of package 165,
specifically the literal 239-coordinate power branch \cite{wide}. This
route uses neither a K215 simulation nor Jones's universal polynomial.
The inherited logical calibration is the named C222-to-ZF refutation
theorem followed by the ZF/ZFC refutation equivalence. The new construction
combines the linear-radius and affine-radix arguments of package 164 with
the self-power horizon and relaxed remainder sphere of package 165.

\subsection{The new selector radius}

The source has the cofinal even horizon and genuine power request
\[
 t=16+2h^2,\qquad F=t^t.
\]
Here $F=t^t$ is the meaning of existing quadratic power certificates;
it is not a variable-exponent operation in the final polynomial.
Change the two source equations
\[
 R=F^2\quad\hbox{to}\quad R=F,
 \qquad q=(2z+1)^2\quad\hbox{to}\quad q=4z+1.
\]
Every power request, frame condition, orientation condition and local
Horn relation is retained. The already relaxed sphere is also retained:
\[
 M^2=2\sum_{j=1}^{11}R_j^2+G(a,b,c),\qquad
 G(a,b,c)=a^2+b^2+c(c+1).
\]
The polynomial $G$ is nonnegative on $\Z^3$ and represents every
nonnegative integer; this standard consequence of the three-square theorem
is also independently formalized in \cite{threesquares}.

\begin{lemma}
If $t\ge16$ and $F=t^t$, then
$F\ge\lcm(1,\ldots,t)^2$.
\end{lemma}
\begin{proof}
Write $\Lambda_n=\lcm(1,\ldots,n)$ and $k=\lceil n/2\rceil$.
Prime valuations in the central binomial coefficient give
$\Lambda_n\mid\Lambda_k\binom nk$.
For $n=2m$, the binomial coefficient is at most $4^m$;
for $n=2m+1$, the two equal middle coefficients together are at most
$2^{2m+1}$, so either is at most $4^m$.
Induction gives $\Lambda_n\le4^n$. Thus
$t^t\ge16^t\ge\Lambda_t^2$.
\end{proof}

This is the exact threshold required by the binomial CRT selector at
$B=\binom{R+t}{t}$. To make both directions explicit, for each
$i\in[1,t]$, the integer $R+i>\Lambda_t^2$ has a prime $p$ with
$v_p(R+i)>2v_p(\Lambda_t)$. Subtracting the valuations at the other
positions gives an exact prime power $p^e\Vert B$ with $p^e>t$ and
$p^e\mid R+i$. Conversely each exact prime power dividing $B$ divides
some $R+i$, by the zero-one summands in Legendre's formula. This supplies
both the local-row extraction and reverse CRT completeness at $R=F$.

\subsection{Signed radix recovery before using the sphere}

The retained beta endpoints and index row are
\[
 y=F(2w+1)^2,\quad x=F+1+(y+1)u_0,\quad
 x=1+[1+(t+1)y]u_t,\quad B\mu=C_0+R+1.
\]
The genuine horizon power gives $F\ge2^{64}$, hence $y\ge F>0$.
The affine equation $q=4z+1$ excludes $q=0,-1$. The retained signed
power theorem authenticates the requests at $q$ and $q+1$. If $q=1$,
the requests at base 2 first make all actual representatives nonnegative.
The literal offset in this branch is
$H=2+\sum$ of the nine actual representatives, so $H\ge2$.
The index row gives $B\ne0$. The self-radix relation $q=(BH)^x$ would
then force $x=0$, contradicting the last beta endpoint. Hence $q\ne1$.

For the remaining signed branches, $|q|\ge3$, the same argument gives
$H\ge2$, $B\ne0$ and $x>0$. The two beta endpoints imply
\[
 x\ge F+1,\qquad x\ge2+(t+1)y,\qquad
 |q|=|BH|^x\ge2^x,\qquad |B|\le |q|/4.
\]
Thus $M=q^x\ne0$ before dividing any remainder row. Normalized rows and
the sphere give integral remainders with
$|r_j|\le|L_j|/\sqrt2<3|L_j|/4$.

For $n=F+t$, all coefficients $\binom nj$, $j\le t$, are bounded by
$(F^2+1)^t$. The endpoint growth makes every genuine low tail smaller
than $|L_j|/32$. The difference between a recovered remainder and its
true tail is a multiple of $L_j$ with magnitude below
$(3/4+1/32)|L_j|<|L_j|$, so it vanishes. In particular the binomial
digit gives $B=\binom{F+t}{t}>0$; hence $q=(BH)^x>0$.
All remaining digits, local rows and labels are then recovered at their
original precision. The complete ordinary argument is retained in
\texttt{direct233/PROOF.md}.

\subsection{The self-power label argument}

The older inequality $(t+1)^t<F$ cannot be used when $F=t^t$.
Instead, suppose equal products of weights $1+jF$, $1\le j\le t$,
have at least one side with at most $t$ factors. Unequal counts would
contradict $F+1>t^d$ for the smaller count $d\le t$. Thus both counts
are the same $d\le t$. Regard the two products as degree-$d$ polynomials
in a new indeterminate. All their coefficients are positive and at most
$\binom dk t^k\le t^d\le F$. Their corresponding differences have
absolute value less than $F$. The lowest nonzero coefficient therefore
cannot vanish after evaluation at $F$, by reduction modulo $F$.
The polynomials agree, and their distinct roots $-1/j$ recover all
label multiplicities. This restores the exact Horn-closure interface.

\subsection{Completeness and the four algebraic projections}

Pad a finite refutation to an allowed even horizon. Because
$v_2(F)\ge t>v_2(t!)$, an arbitrarily large odd factor can make the
existing beta spacing $y=F(2w+1)^2$ an even multiple of $t!$.
The beta moduli are then pairwise coprime and odd; their code $x$ can
be chosen even. Also
\[
 v_2\binom{F+t}{t}=0,
\]
since $v_2(F+i)=v_2(i)$ for $1\le i\le t$. The retained CRT
representative lifts can therefore make $H$ odd while preserving
dominance and parity. The resulting $q=(BH)^x$ is an odd square and
$z=(q-1)/4$ is integral. The canonical remainders are below $M/32$,
so the nonnegative sphere gap has a $G$ representation. This constructs
a zero of the edited source from every refutation.

The two edited rows are monic affine equations. Delete them after their
exact substitutions to obtain 237 variables. Three directions in the
common kernel of the residual Hessians permit integral polynomial
elimination without increasing degree. A further affine label-definition
row removes one coordinate, giving 233 variables and 169 residuals.
Each projection has its recorded integral inverse; every row-basis
operation has its elementary inverse. These transformations preserve
integer zero sets, though they need not preserve the value of the sum
of squares away from zeros. The final polynomial is the literal sum
of the new squared residuals, and all its coefficients are included.

\section{What was verified and what is inherited}

The release contains actual coefficient payloads, complete equation graphs
or residual lists, proofs, provenance, and standard-library checkers. It
does not require network access to reconstruct its new finite artifacts.
The following are different kinds of evidence and are kept separate.

\begin{enumerate}
\item \textbf{Ordinary unbounded arguments.} These prove the new alignment,
carry, controller, signed-Pell, selector, and completeness claims for
arbitrary witnesses. The proofs compose with specifically identified
native or C222 source calibrations. All are ordinary finitary arithmetic
arguments formalizable in ZFC; no standard-witness-only inference is used.
\item \textbf{Exact finite checks.} These reconstruct the full constants,
authenticate their bytes, check all support identities and controller
conditions, verify graph liveness and exact top homogeneous parts, and
verify every new projection, inverse row operation and expanded quartic
coefficient. Independent reviews found and checked the actual stored-to-
semantic representative offsets and the order of denominator recovery.
\item \textbf{Inherited logical proofs.} K215/native calibration is used
only on the two compact branches; the direct branch uses the existing
C222, ZF and ZFC refutation theorems. Those entire historical developments
were not rebuilt here. Some historical dependency files advertised by
package 165 are absent from the inspected tree; no unavailable old checker
is claimed to have run. The new finite source-interface commutation and
all new algebra were checked from the supplied literal source.
\end{enumerate}

The distinction is substantive. A circuit checker cannot prove an
unbounded Pell classification, and an exact polynomial identity cannot
supply a missing positivity theorem. Conversely, inheriting a valid source
calibration does not prevent a new local transformation from being a
complete mathematical proof when its interfaces and both directions are
checked. This is the status of the constructions in this article.

\paragraph{A rejected transfer.}
The trace substitution in Section~\ref{sec:85} uses the literal identity
$f\equiv1\pmod c$. P66 has a free strong Pell root for which only
$f^2\equiv1\pmod c$ is immediate. That does not imply $f\equiv1\pmod c$.
The archive includes a general Pell-family obstruction to this proposed
shortcut. We do not claim that simply replacing P66's $f^2(cT-R)-c$
by $f(cT-R)-c$ is sound.

\section{Further research questions}

The improvements suggest several concrete next problems.
\begin{enumerate}
\item Can the Start-only layout be shortened using a larger authenticated
label pattern, while retaining a transparent all-offset alignment proof?
The new layout already uses one specific tile to eliminate the last
false offset; a general search must certify semantic patterns, not merely
physical support.
\item Can square transport's carrier-degree reduction from 57 to 21
lower the number of quadratic witnesses in a direct compilation, even
though it leaves the total degree 232 unchanged?
\item Can the low-variable trace form be specialized to the new complete
Start-only native tuple with a separately documented full interface?
The present ten-variable release deliberately retains its authenticated
radix-67 tuple; the coherent new specialization is not silently assumed.
\item Can the direct selector radius be lowered below $t^t$ while preserving
both prime-power coverage and exact recovery of label multiplicities?
Any proposed bound must establish the two properties together.
\item Can additional nonlinear, integrally invertible projections reduce
the direct 233-coordinate system? The current common-radical and affine
certificates give exact local reductions, not a lower bound on all
possible transformations.
\item An end-to-end formal derivation inside the repository's object-ZFC
calculus would strengthen the deliverable's verification status. Its
remaining work includes the complete native or C222 arithmetic calibration,
not merely rechecking the finite JSON arrays.
\end{enumerate}

\appendix
\section{Complete fixed language for the compact equation}
\label{app:language}
\providecommand{\fl}[1]{\left\lfloor#1\right\rfloor}
\providecommand{\rg}[2]{\{#1,\ldots,#2\}}
Every item in this appendix defines a fixed constant. All index ranges are finite.
Use the actions $u_j,v_j,d_j$ decoded below. For $d=0,1$, let
\[
\begin{gathered}
h(q,b)=2q+b+2,\quad
\mathcal S_d=\{u_j:d_j=d,\ 0\le j<430\},\\
r_{d,q}=205-199d+|\mathcal S_d\cap\rg0{q-1}|.
\end{gathered}
\]
An interval ending at $-1$ is empty. The table defines $t_i$ for
$0\le i\le942$; all index ranges are fixed:
\[
\setlength{\arraycolsep}{1.4pt}\begin{array}{c|l}
\text{index }i&\text{tile }t_i=(w_i,n_i,p_i,s_i)\\\hline
0,1&(2,435,3,1),\ (5,0,0,0)\\
j+2\ (j=2q+b<430)&((1-d_j)r_{0,u_j},h(q,b),d_jr_{1,u_j},v_j)\\
434-2a+b\ (a,b=0,1)&(0,432a+b,0,432a+b)\\
\substack{424+2r_{d,q}+b\\q\in\mathcal S_d,\ d,b=0,1}&(dr_{d,q},b,(1-d)r_{d,q},h(q,b))\\
924+2a+b\ (a,b=0,1)&(0,432a+b,0,438)\\
928&(2,434,2,0)\\
929&(3,436,4,4)\\
930&(4,436,4,0)\\
497+t\ (t=434,435,436)&(1,438,1,t)\\
934+3a+b\ (a,b=0,1,2)&(a^2,437,b,437)
\end{array}
\]
\textbf{A closed decoder for the 430 actions.}
Put $D_b(A,k)=\fl{A/b^k}\bmod b$, where
$a\bmod b=a-b\fl{a/b}$. For $0\le q<215$, $b=0,1$, $j=2q+b$:
\[
\begin{gathered}
N=\mathcal A\bmod2^{636},\quad
T_0=\fl{\mathcal A/2^{636}}\bmod2^{1675},\\
E_0=\fl{\mathcal A/2^{2311}},\\
\nu_j=D_2(N,j),\quad \sigma_j=\sum_{k=0}^{j-1}\nu_k,\\
m_q=D_2(N,430+q),\quad \tau_q=\sum_{k=0}^{q-1}m_k,\\
F_q=14+m_q(D_{16}(E_0,\tau_q)-14),\\
u_j=\nu_j(1+\sigma_j)+(1-\nu_j)D_{216}(T_0,j-\sigma_j),\\
v_j=D_2(F_q,2b),\quad d_j=D_2(F_q,2b+1).
\end{gathered}
\]
Empty sums are zero. Here $(v_0,d_0,u_0)=(1,1,1)$; the constructor tile
uses this actual first target, including its state renaming.

\begin{samepage}
\textbf{The integer $\mathcal A$.}
The alphabet, in order of digit values $0,\ldots,61$, is\par
\noindent
{\small\ttfamily 0123456789ABCDEFGHIJKLMNOPQRSTUVWXYZabcdefghijklmnopqrstuvwxyz}.\par
Concatenate the following 436 characters as digits $e_0,\ldots,e_{435}$;
put $\mathcal A=\sum_{j=0}^{435}e_j62^{435-j}$.\par
\begingroup\small\ttfamily
BxEnhvy6NipzkWmROo0GsA3uK3twWLxscLq20UTumy5Cdu461LwD6g7RDNzF8Y\par
iYDFTpB8ehAypd7d6OsOL4APamnJeNXQypUezCmW9n4oXP6XfgmiWuAWZ8ZJ7U\par
OP7gu4kVdlGZqZHULQMar1dXNeJW4i7oNVI8usp8BlovjjZbwPCovZI328fRUq\par
4q92R81s5GjWQaQfGhkq9UYMK3rRMaBmsOIBkJqc2cj2VRIZ6mjzkBIOkW9SjZ\par
FzeA6J5u4j3sBL3SdgwXzM9mHvQivu8EJ1TBTpRhHDqZIbmYgwQ8VqTAFotl6F\par
We1YC554GKrVHsA8oUGKSA6F6k5KREBSD0Tgd7bFALY84LLbpw574BQyvdoHrl\par
0Oe35uwDf15DpT6VJWZBjSbHPpPAINYTte0eUxnTFTKBlGj8Am6sb1gLldh8WZ\par
8P\par
\endgroup
\end{samepage}


\section{Complete coefficient recipe for the ten-variable equation}
\label{app:85constants}
% Exact finite recipe extracted from the retained indexed_constants.tex.
% In this appendix K denotes the fixed port K_0, and all other symbols are constants.
\begingroup
\providecommand{\fl}[1]{\left\lfloor#1\right\rfloor}
\providecommand{\rg}[2]{\{#1,\ldots,#2\}}
In this appendix $K$ denotes the port $K_0$ in the new equations.
Every symbol introduced here is a fixed constant, independent of the ten
quantified variables. The formulas and 436-character seed are the retained
indexed BFS943 radix67 specification; they reconstruct the exact supplied
payloads, with $\beta=2^{590069}-1$.

\textbf{The four fixed integers.}
Every index, set, sum, product and floor below defines a constant, not
another unknown. Put
\[
\rho=2^{67},\quad B=\rho^{8807},\quad\beta=B-1,\quad
G_m(x)=\sum_{j=0}^{m-1}x^j.
\]
Define
\[
\begin{split}
\mathcal E={}&\rg0{942}\cup\rg{983}{992}\cup\rg{1003}{1011}\\
&\cup\rg{1013}{1031}\cup\{1063,1064,1095,2195,3296\}.
\end{split}
\]
Then
\begin{align*}
M_C={}&B-1-\sum_{e\in\mathcal E\setminus\{1\}}\rho^e-2\rho^{1063},\\
M_F={}&B-1+4+(2^{51}-2)(\rho^{4393}+\rho^{5495})\\
&+\rho^{1013}G_{10}(\rho)+\rho^{1065}G_9(\rho)
  +\rho^{2195}+\rho^{3296}\\
&+(2^{65}-1)\rho^{8792}G_{13}(\rho)+(2^{19}-1)\rho^{8805},\\
K={}&\rho^{30}+\rho^{62}+\rho^{1100}+\rho^{2201}+\rho^{8849}\\
&+(\rho^{4393}+\rho^{5495})\Phi(1/\rho),
\end{align*}
where
\[
\begin{split}
\Phi(x)={}&(1+2^{51})G_{943}(x)
 -2^{10}\sum_{i=0}^{942}v(t_i)x^i\\
&+\sum_{(b,e,m)\in\mathcal C}2^b x^e G_m(2x),\\
\mathcal C={}&\{(10,1023,8),(18,1013,9),(27,1003,8),\\
&(35,983,9),(44,1031,1),(45,1022,1),(46,1011,1),\\
&(47,992,1),(48,1095,1),(49,2195,1),(50,3296,1)\}.
\end{split}
\]
The rational intermediate in K cancels exactly; all four results are positive integers.

\textbf{The 943 indexed fixed tiles.}
Use the actions $u_j,v_j,d_j$ decoded below. Let
\[
\begin{gathered}
h(q,b)=2(q+1)+b,\quad
\mathcal R=\{u_j:d_j=1,\ 0\le j<430\},\\
\mathcal L=\{u_j:d_j=0,\ 0\le j<430\},\\
r_q=6+|\mathcal R\cap\rg0{q-1}|,\quad
l_q=205+|\mathcal L\cap\rg0{q-1}|.
\end{gathered}
\]
An interval ending at $-1$ is empty. The table defines $t_i$ for
$0\le i\le942$; all index ranges are fixed:
\[
\setlength{\arraycolsep}{1.4pt}\begin{array}{c|l}
\text{index }i&\text{tile }t_i=(w,n,e,s)\\\hline
0,1&(2,435,3,1),\ (5,0,0,0)\\
j+2\ (j=2q+b<430)&((1-d_j)l_{u_j},h(q,b),d_jr_{u_j},v_j)\\
432+b\ (b=0,1)&(0,h(215,b),0,h(215,b))\\
434+b\ (b=0,1)&(0,b,0,b)\\
424+2r_q+b\ (q\in\mathcal R,b=0,1)&(r_q,b,0,h(q,b))\\
424+2l_q+b\ (q\in\mathcal L,b=0,1)&(0,b,l_q,h(q,b))\\
924+j\ (0\le j<4)&(0,j+430\fl{j/2},0,438)\\
928&(2,434,2,0)\\
929&(3,436,4,4)\\
930&(4,436,4,0)\\
497+t\ (t=434,435,436)&(1,438,1,t)\\
\substack{934+3a+b-6\fl{a/4}\\a\in\{0,1,4\},\ b=0,1,2}&(a,437,b,437)
\end{array}
\]
For $t=(w,n,e,s)$ put
\[
v(t)=w+2^8n+2^{17}e+2^{25}s+127\,2^{34}.
\]

\textbf{A closed decoder for the 430 actions.}
Put $D_b(A,k)=\fl{A/b^k}\bmod b$, where
$a\bmod b=a-b\fl{a/b}$. For $0\le q<215$, $b=0,1$, $j=2q+b$:
\[
\begin{gathered}
N=\mathcal A\bmod2^{636},\quad
T_0=\fl{\mathcal A/2^{636}}\bmod2^{1675},\\
E_0=\fl{\mathcal A/2^{2311}},\\
\nu_j=D_2(N,j),\quad \sigma_j=\sum_{k=0}^{j-1}\nu_k,\\
m_q=D_2(N,430+q),\quad \tau_q=\sum_{k=0}^{q-1}m_k,\\
F_q=14(1-m_q)+m_qD_{16}(E_0,\tau_q),\\
u_j=\nu_j(1+\sigma_j)+(1-\nu_j)D_{216}(T_0,j-\sigma_j),\\
v_j=D_2(F_q,2b),\quad d_j=D_2(F_q,2b+1).
\end{gathered}
\]
Empty sums are zero. Here $(v_0,d_0,u_0)=(1,1,1)$; the constructor tile
uses this actual first target, including its state renaming.

\begin{samepage}
\textbf{The integer $\mathcal A$.}
The alphabet, in order of digit values $0,\ldots,61$, is\par
\noindent
{\small\ttfamily 0123456789ABCDEFGHIJKLMNOPQRSTUVWXYZabcdefghijklmnopqrstuvwxyz}.\par
Concatenate the following 436 characters as digits $e_0,\ldots,e_{435}$;
put $\mathcal A=\sum_{j=0}^{435}e_j62^{435-j}$.\par
\begingroup\small\ttfamily
BxEnhvy6NipzkWmROo0GsA3uK3twWLxscLq20UTumy5Cdu461LwD6g7RDNzF8Y\par
iYDFTpB8ehAypd7d6OsOL4APamnJeNXQypUezCmW9n4oXP6XfgmiWuAWZ8ZJ7U\par
OP7gu4kVdlGZqZHULQMar1dXNeJW4i7oNVI8usp8BlovjjZbwPCovZI328fRUq\par
4q92R81s5GjWQaQfGhkq9UYMK3rRMaBmsOIBkJqc2cj2VRIZ6mjzkBIOkW9SjZ\par
FzeA6J5u4j3sBL3SdgwXzM9mHvQivu8EJ1TBTpRhHDqZIbmYgwQ8VqTAFotl6F\par
We1YC554GKrVHsA8oUGKSA6F6k5KREBSD0Tgd7bFALY84LLbpw574BQyvdoHrl\par
0Oe35uwDf15DpT6VJWZBjSbHPpPAINYTte0eUxnTFTKBlGj8Am6sb1gLldh8WZ\par
8P\par
\endgroup
\end{samepage}

\endgroup



\section{Every quadratic residual of the balanced direct equation}
\label{app:residuals}
The following list is generated directly from the verified sparse arrays.
All coordinates $x_0,\ldots,x_{232}$ range over $\Z$. The indexing of the
printed residuals begins at zero; summing their squares gives
\eqref{eq:direct}. Numerical subscripts are part of variable names, not
extra parameters.
\begingroup\footnotesize
\setlength{\jot}{1pt}
\[\begin{aligned}
r_{0}={}&x_{0}-x_{0}x_{2}-x_{2}+x_{7}
\end{aligned}\]
\[\begin{aligned}
r_{1}={}&-x_{2}x_{5}-x_{3}+x_{6}x_{7}
\end{aligned}\]
\[\begin{aligned}
r_{2}={}&-1+x_{1}-x_{6}x_{7}
\end{aligned}\]
\[\begin{aligned}
r_{3}={}&1+x_{3}-x_{4}x_{18}+x_{16}
\end{aligned}\]
\[\begin{aligned}
r_{4}={}&-1+x_{8}-x_{8}^{2}-2x_{8}x_{9}-x_{9}-x_{9}^{2}+x_{19}
\end{aligned}\]
\[\begin{aligned}
r_{5}={}&1+x_{1}+x_{2}-x_{3}x_{9}-x_{3}x_{231}-x_{6}x_{7}-12x_{8}-10x_{8}^{2}-6x_{8}x_{9}+2x_{8}x_{10}
\\ &-16x_{8}x_{19}-2x_{9}-2x_{9}x_{10}-8x_{9}x_{19}+4x_{10}-x_{10}^{2}-2x_{10}x_{11}-x_{11}^{2}-11x_{19}-4x_{19}^{2}
\\ &+x_{20}+2x_{22}+3x_{23}+2x_{24}+2x_{26}+2x_{28}+2x_{30}+4x_{31}+4x_{31}^{2}+x_{32}+4x_{34}+2x_{35}
\\ &-4x_{36}+4x_{36}^{2}-x_{37}+2x_{38}+x_{39}+2x_{41}-4x_{60}^{2}+3x_{61}+x_{75}+4x_{202}x_{216}-x_{210}
\end{aligned}\]
\[\begin{aligned}
r_{6}={}&-1-x_{1}+x_{6}x_{7}-x_{8}-x_{8}^{2}-4x_{8}x_{20}-4x_{20}-4x_{20}^{2}+x_{21}
\end{aligned}\]
\[\begin{aligned}
r_{7}={}&1+2x_{21}-4x_{21}^{2}+x_{22}
\end{aligned}\]
\[\begin{aligned}
r_{8}={}&-1-x_{1}+x_{6}x_{7}-x_{9}^{2}-2x_{9}x_{10}-x_{10}-x_{10}^{2}+x_{23}
\end{aligned}\]
\[\begin{aligned}
r_{9}={}&-1+x_{8}-x_{8}^{2}-2x_{8}x_{10}-x_{9}-x_{10}^{2}+x_{24}
\end{aligned}\]
\[\begin{aligned}
r_{10}={}&2+2x_{0}+x_{1}+x_{2}-x_{3}x_{8}-x_{3}x_{9}-x_{3}x_{231}+x_{4}x_{18}-x_{6}x_{7}-6x_{8}-10x_{8}^{2}
\\ &-9x_{8}x_{9}-17x_{8}x_{10}+4x_{8}x_{23}-24x_{8}x_{24}+12x_{9}-x_{9}^{2}-3x_{9}x_{10}-8x_{9}x_{19}+16x_{9}x_{23}
\\ &+8x_{9}x_{24}+4x_{9}x_{47}-7x_{10}+10x_{10}^{2}-2x_{10}x_{11}-20x_{10}x_{19}-8x_{10}x_{24}+8x_{10}x_{37}
\\ &-x_{11}^{2}-x_{16}-4x_{19}+8x_{19}^{2}-8x_{19}x_{24}+2x_{22}+5x_{23}-8x_{23}x_{24}-12x_{24}-8x_{24}^{2}+x_{25}
\\ &+2x_{26}+x_{27}+2x_{28}+2x_{30}+12x_{31}+12x_{31}^{2}-x_{32}+4x_{34}+2x_{35}+3x_{36}-4x_{36}^{2}-3x_{37}
\\ &+4x_{37}^{2}+x_{38}+x_{39}+2x_{41}-2x_{43}+4x_{47}^{2}-x_{48}+4x_{54}-4x_{54}^{2}+x_{55}-7x_{58}+8x_{58}^{2}
\\ &+4x_{59}^{2}+3x_{60}+2x_{61}+x_{71}+x_{75}+x_{76}+x_{136}x_{190}+x_{190}-x_{191}+4x_{202}x_{216}
\\ &-4x_{206}^{2}-x_{210}
\end{aligned}\]
\[\begin{aligned}
r_{11}={}&-1+6x_{8}+x_{8}x_{9}+3x_{8}x_{10}-8x_{8}x_{23}+3x_{9}+6x_{9}^{2}+3x_{9}x_{10}+2x_{10}-10x_{10}^{2}
\\ &+2x_{10}x_{11}+24x_{10}x_{25}+x_{11}^{2}+2x_{19}+4x_{23}-6x_{24}-4x_{25}^{2}+x_{26}-4x_{31}-4x_{31}^{2}+x_{32}
\\ &-x_{39}-2x_{43}-13x_{71}
\end{aligned}\]
\[\begin{aligned}
r_{12}={}&-1+x_{0}-x_{1}+x_{3}x_{8}+x_{3}x_{9}-x_{4}x_{18}+x_{6}x_{7}-7x_{8}+x_{8}^{2}-5x_{8}x_{9}-x_{8}x_{10}
\\ &+8x_{8}x_{23}+8x_{8}x_{24}-2x_{9}+2x_{9}^{2}+x_{9}x_{10}+8x_{9}x_{23}+8x_{9}x_{24}+2x_{10}^{2}+8x_{10}x_{23}
\\ &+8x_{10}x_{24}+x_{16}-3x_{19}+4x_{23}-4x_{23}^{2}-8x_{23}x_{24}+5x_{24}-4x_{24}^{2}+x_{27}+4x_{31}+4x_{31}^{2}
\\ &-x_{32}+x_{71}-x_{75}-x_{76}-x_{136}x_{190}-x_{190}+x_{191}-2x_{206}^{2}
\end{aligned}\]
\[\begin{aligned}
r_{13}={}&1+x_{1}+x_{2}-x_{3}x_{9}-x_{3}x_{231}-x_{6}x_{7}-3x_{8}-x_{8}^{2}-3x_{8}x_{9}+x_{8}x_{10}+2x_{9}-3x_{9}^{2}
\\ &-3x_{9}x_{10}-8x_{9}x_{27}+5x_{10}-3x_{10}^{2}-8x_{10}x_{27}-x_{19}-2x_{23}+4x_{27}-4x_{27}^{2}+x_{28}+x_{71}
\\ &+x_{75}
\end{aligned}\]
\[\begin{aligned}
r_{14}={}&2+3x_{0}+x_{0}x_{2}+2x_{1}+x_{2}-x_{2}x_{5}-x_{3}-x_{3}x_{8}-x_{3}x_{9}+x_{4}x_{18}-x_{6}x_{7}-x_{7}
\\ &-10x_{8}+2x_{8}^{2}-5x_{8}x_{9}-11x_{8}x_{10}+4x_{8}x_{23}-8x_{8}x_{24}-4x_{8}x_{49}-4x_{9}^{2}-5x_{9}x_{10}
\\ &-16x_{9}x_{23}+12x_{9}x_{24}+8x_{9}x_{27}-11x_{10}-2x_{10}^{2}+4x_{10}x_{11}-16x_{10}x_{23}-8x_{10}x_{24}
\\ &+8x_{10}x_{27}-2x_{11}+x_{11}^{2}-x_{16}+x_{19}+2x_{22}-25x_{23}-8x_{23}^{2}-8x_{23}x_{24}-13x_{24}-8x_{24}^{2}
\\ &+2x_{26}-3x_{27}+4x_{27}^{2}+x_{28}+x_{29}+2x_{30}+12x_{31}+12x_{31}^{2}-x_{32}+4x_{34}+2x_{35}-12x_{36}
\\ &+12x_{36}^{2}-3x_{37}+2x_{38}-2x_{39}+2x_{41}+12x_{49}-4x_{49}^{2}+x_{50}-4x_{54}+4x_{54}^{2}-x_{55}-4x_{60}^{2}
\\ &+3x_{61}+x_{71}+x_{71}^{2}-x_{72}+x_{75}+x_{76}+4x_{202}x_{216}-8x_{206}^{2}-x_{210}
\end{aligned}\]
\[\begin{aligned}
r_{15}={}&1+x_{0}+x_{8}+7x_{8}^{2}-8x_{8}x_{9}+4x_{8}x_{10}+4x_{8}x_{23}-3x_{9}-4x_{9}x_{10}-8x_{9}x_{19}
\\ &-8x_{9}x_{24}+8x_{9}x_{27}+16x_{10}-9x_{10}^{2}-2x_{10}x_{11}-20x_{10}x_{19}+16x_{10}x_{24}-24x_{10}x_{25}
\\ &+8x_{10}x_{27}-32x_{10}x_{29}+4x_{10}x_{57}-x_{11}^{2}-7x_{19}+8x_{19}^{2}-8x_{19}x_{24}-3x_{23}+6x_{24}
\\ &+4x_{24}^{2}+x_{25}+4x_{25}^{2}-x_{26}-4x_{27}+4x_{27}^{2}-x_{28}-2x_{29}-4x_{29}^{2}+x_{30}-8x_{31}-8x_{31}^{2}
\\ &+2x_{32}+7x_{36}-8x_{36}^{2}+2x_{37}+x_{39}+x_{43}-4x_{54}+4x_{54}^{2}-x_{55}+4x_{57}^{2}-3x_{58}+4x_{58}^{2}
\\ &-x_{59}-2x_{206}^{2}
\end{aligned}\]
\[\begin{aligned}
r_{16}={}&4x_{8}^{2}+10x_{8}x_{9}+6x_{8}x_{10}-8x_{8}x_{23}-8x_{8}x_{24}-14x_{9}+2x_{9}^{2}+8x_{9}x_{10}
\\ &+8x_{9}x_{23}-8x_{9}x_{24}-4x_{9}x_{47}+2x_{10}-5x_{10}^{2}+2x_{10}x_{11}+4x_{10}x_{19}-8x_{10}x_{24}+x_{11}^{2}
\\ &+11x_{19}-8x_{19}x_{23}-2x_{22}+16x_{23}-4x_{23}^{2}+8x_{23}x_{24}-4x_{24}+4x_{24}^{2}-2x_{26}-x_{27}
\\ &-2x_{28}-2x_{30}-3x_{31}-4x_{31}^{2}-x_{32}-4x_{34}-2x_{35}-8x_{36}+8x_{36}^{2}-2x_{37}-2x_{38}-x_{39}
\\ &-2x_{41}+x_{43}-4x_{47}^{2}+x_{48}+7x_{58}-8x_{58}^{2}-4x_{59}^{2}-3x_{60}-2x_{61}+4x_{71}-4x_{202}x_{216}
\\ &+x_{210}
\end{aligned}\]
\[\begin{aligned}
r_{17}={}&-1-x_{1}+x_{6}x_{7}-5x_{8}-x_{8}^{2}-2x_{8}x_{10}+x_{9}+x_{9}^{2}+2x_{9}x_{10}-x_{10}-x_{23}-x_{24}
\\ &-4x_{31}-4x_{31}^{2}+x_{32}
\end{aligned}\]
\[\begin{aligned}
r_{18}={}&1+x_{1}+x_{2}-x_{3}x_{8}-x_{3}x_{9}-x_{3}x_{231}+x_{4}x_{18}-x_{6}x_{7}+10x_{8}-6x_{8}^{2}+4x_{8}x_{9}
\\ &-2x_{8}x_{10}-32x_{8}x_{19}+36x_{8}x_{20}+9x_{9}+4x_{9}^{2}-4x_{9}x_{10}-24x_{9}x_{19}+32x_{9}x_{20}-4x_{10}
\\ &+x_{10}^{2}-8x_{10}x_{37}-x_{16}+6x_{19}-12x_{19}^{2}+32x_{19}x_{20}-9x_{20}-28x_{20}^{2}+4x_{21}+4x_{21}^{2}
\\ &-x_{22}-4x_{23}-7x_{24}-4x_{31}-4x_{31}^{2}+x_{32}+x_{33}+8x_{36}-8x_{36}^{2}+6x_{37}-4x_{37}^{2}+x_{38}
\\ &-8x_{58}+8x_{58}^{2}+4x_{59}^{2}-x_{60}-4x_{71}+x_{75}+x_{76}+x_{136}x_{190}+x_{190}-x_{191}+x_{207}^{2}
\\ &+x_{208}^{2}-x_{217}^{2}
\end{aligned}\]
\[\begin{aligned}
r_{19}={}&1+2x_{33}-4x_{33}^{2}+x_{34}
\end{aligned}\]
\[\begin{aligned}
r_{20}={}&-x_{9}^{2}-4x_{9}x_{34}-4x_{34}^{2}+x_{35}
\end{aligned}\]
\[\begin{aligned}
r_{21}={}&-5x_{8}-6x_{8}^{2}-2x_{8}x_{9}-2x_{8}x_{10}-16x_{8}x_{19}-16x_{8}x_{27}+2x_{9}+6x_{9}^{2}-6x_{9}x_{10}
\\ &+8x_{9}x_{19}-8x_{9}x_{27}+4x_{10}^{2}-8x_{10}x_{19}+4x_{10}x_{27}+x_{19}-8x_{19}^{2}-8x_{19}x_{27}-x_{23}
\\ &-3x_{24}+4x_{27}+x_{28}+4x_{31}+4x_{31}^{2}-x_{32}-3x_{36}+4x_{36}^{2}-x_{37}-x_{42}-x_{58}+10x_{71}
\end{aligned}\]
\[\begin{aligned}
r_{22}={}&-1-x_{1}+x_{6}x_{7}-3x_{8}-7x_{9}+x_{9}^{2}+2x_{9}x_{10}-3x_{10}+x_{10}^{2}-3x_{23}-x_{24}+4x_{36}
\\ &-4x_{36}^{2}+x_{37}
\end{aligned}\]
\[\begin{aligned}
r_{23}={}&-1-5x_{8}-x_{8}^{2}-x_{8}x_{9}-x_{8}x_{10}-x_{9}^{2}-x_{9}x_{10}+3x_{10}-3x_{10}^{2}-8x_{10}x_{37}+x_{19}
\\ &-2x_{24}+4x_{37}-4x_{37}^{2}+x_{38}-x_{71}
\end{aligned}\]
\[\begin{aligned}
r_{24}={}&-x_{10}-x_{10}^{2}-2x_{10}x_{11}-x_{11}^{2}+x_{39}
\end{aligned}\]
\[\begin{aligned}
r_{25}={}&1+2x_{8}-x_{8}^{2}+2x_{8}x_{9}+6x_{8}x_{10}-8x_{8}x_{23}-8x_{8}x_{24}-4x_{8}x_{39}+6x_{9}^{2}
\\ &+2x_{9}x_{10}-8x_{9}x_{24}-8x_{9}x_{27}-11x_{10}+x_{10}^{2}+4x_{10}x_{11}+8x_{10}x_{23}+16x_{10}x_{24}
\\ &-8x_{10}x_{27}-8x_{10}x_{37}-4x_{10}x_{57}+2x_{11}^{2}+4x_{19}-2x_{22}-3x_{23}+4x_{23}^{2}+8x_{24}-2x_{26}
\\ &+3x_{27}-4x_{27}^{2}-x_{28}-2x_{30}-2x_{32}-4x_{34}-2x_{35}-4x_{36}+4x_{36}^{2}+3x_{37}-4x_{37}^{2}-x_{38}
\\ &-4x_{39}^{2}+x_{40}-2x_{41}-x_{43}+4x_{54}-4x_{54}^{2}+x_{55}-4x_{57}^{2}+4x_{58}-4x_{58}^{2}+x_{59}-4x_{60}
\\ &-2x_{61}-16x_{71}-4x_{202}x_{216}+x_{210}
\end{aligned}\]
\[\begin{aligned}
r_{26}={}&1+x_{1}+x_{2}-x_{3}x_{9}-x_{3}x_{231}-x_{6}x_{7}+2x_{8}+2x_{8}x_{9}+2x_{8}x_{10}-8x_{8}x_{39}+x_{9}^{2}
\\ &+5x_{10}-9x_{10}^{2}+4x_{10}x_{11}-16x_{10}x_{40}+2x_{11}^{2}-x_{23}-2x_{24}+2x_{39}-8x_{39}^{2}-4x_{40}^{2}
\\ &+x_{41}-2x_{71}+x_{75}
\end{aligned}\]
\[\begin{aligned}
r_{27}={}&-1-x_{1}+x_{6}x_{7}+3x_{8}-15x_{8}^{2}+10x_{8}x_{9}+16x_{8}x_{10}-16x_{8}x_{19}+16x_{8}x_{27}-5x_{9}
\\ &+6x_{9}^{2}+2x_{9}x_{10}+8x_{9}x_{27}-4x_{10}^{2}+12x_{10}x_{19}-4x_{10}x_{27}-x_{19}-4x_{19}^{2}+8x_{19}x_{27}
\\ &+x_{23}+2x_{24}-4x_{27}-x_{28}+x_{42}-4x_{71}
\end{aligned}\]
\[\begin{aligned}
r_{28}={}&-1+2x_{8}x_{9}-2x_{8}x_{10}+2x_{9}-5x_{9}^{2}+4x_{9}x_{10}-8x_{9}x_{23}-2x_{10}+x_{10}^{2}
\\ &+8x_{10}x_{23}+x_{19}+6x_{23}-4x_{23}^{2}-x_{24}+x_{43}-4x_{71}
\end{aligned}\]
\[\begin{aligned}
r_{29}={}&-1-x_{10}-x_{10}^{2}-4x_{10}x_{43}-4x_{43}^{2}+x_{44}
\end{aligned}\]
\[\begin{aligned}
r_{30}={}&1+x_{9}-x_{9}^{2}-4x_{9}x_{44}-4x_{44}^{2}+x_{45}
\end{aligned}\]
\[\begin{aligned}
r_{31}={}&-1-6x_{45}-4x_{45}^{2}+x_{46}
\end{aligned}\]
\[\begin{aligned}
r_{32}={}&-1-x_{10}-x_{10}^{2}-4x_{10}x_{46}-4x_{46}^{2}+x_{47}
\end{aligned}\]
\[\begin{aligned}
r_{33}={}&1+x_{1}-x_{6}x_{7}+x_{9}-x_{9}^{2}-4x_{9}x_{47}-4x_{47}^{2}+x_{48}
\end{aligned}\]
\[\begin{aligned}
r_{34}={}&-2x_{42}-4x_{42}^{2}-8x_{42}x_{48}-4x_{48}^{2}+x_{49}
\end{aligned}\]
\[\begin{aligned}
r_{35}={}&-2-2x_{1}-x_{2}+x_{2}x_{5}+x_{3}x_{8}+x_{3}x_{9}+x_{3}x_{231}+x_{6}x_{7}+5x_{8}+x_{8}^{2}+2x_{8}x_{9}
\\ &+2x_{8}x_{10}-4x_{8}x_{49}+2x_{9}+x_{9}^{2}+x_{10}^{2}-x_{19}-x_{24}+12x_{49}-4x_{49}^{2}+x_{50}-x_{75}-x_{76}
\end{aligned}\]
\[\begin{aligned}
r_{36}={}&-1-6x_{50}-4x_{50}^{2}+x_{51}
\end{aligned}\]
\[\begin{aligned}
r_{37}={}&-x_{10}-x_{10}^{2}-4x_{10}x_{51}-4x_{51}^{2}+x_{52}
\end{aligned}\]
\[\begin{aligned}
r_{38}={}&-x_{8}-x_{8}^{2}-4x_{8}x_{52}-4x_{52}-4x_{52}^{2}+x_{53}
\end{aligned}\]
\[\begin{aligned}
r_{39}={}&1-x_{0}+x_{1}-x_{2}x_{5}-x_{3}x_{8}-x_{3}x_{9}-x_{9}-x_{9}^{2}-2x_{9}x_{10}-x_{10}^{2}-14x_{53}-4x_{53}^{2}
\\ &+x_{54}+x_{75}+x_{76}+2x_{206}^{2}
\end{aligned}\]
\[\begin{aligned}
r_{40}={}&1-x_{8}x_{9}+x_{8}x_{10}-x_{9}+x_{9}x_{10}+x_{10}-x_{10}^{2}-2x_{23}+4x_{54}-4x_{54}^{2}+x_{55}+x_{71}
\end{aligned}\]
\[\begin{aligned}
r_{41}={}&2-x_{0}+x_{1}+x_{2}-x_{2}x_{5}-x_{3}x_{8}-x_{3}x_{9}-x_{3}x_{231}+x_{9}-x_{9}^{2}-4x_{9}x_{55}-14x_{48}
\\ &-4x_{48}^{2}-4x_{55}^{2}+x_{56}+x_{75}+x_{76}+2x_{206}^{2}
\end{aligned}\]
\[\begin{aligned}
r_{42}={}&1+2x_{9}-2x_{9}^{2}-8x_{9}x_{55}-8x_{55}^{2}-4x_{56}^{2}+x_{57}
\end{aligned}\]
\[\begin{aligned}
r_{43}={}&-x_{8}+x_{8}^{2}+5x_{8}x_{9}-3x_{8}x_{10}+3x_{9}-25x_{9}^{2}+7x_{9}x_{10}+24x_{9}x_{19}+8x_{9}x_{23}
\\ &+4x_{9}x_{34}+4x_{9}x_{44}+4x_{9}x_{47}+3x_{10}+x_{10}^{2}-2x_{10}x_{11}-4x_{10}x_{19}-8x_{10}x_{23}
\\ &+4x_{10}x_{57}-x_{11}^{2}-6x_{19}-4x_{19}^{2}-7x_{23}+4x_{23}^{2}+2x_{24}+4x_{34}^{2}-x_{35}+x_{39}-x_{43}
\\ &+4x_{44}^{2}-x_{45}+4x_{47}^{2}-x_{48}+4x_{57}^{2}-3x_{58}+4x_{58}^{2}-x_{59}+x_{71}
\end{aligned}\]
\[\begin{aligned}
r_{44}={}&x_{8}-x_{8}^{2}-x_{8}x_{9}-x_{8}x_{10}+4x_{9}-x_{9}^{2}-x_{9}x_{10}+2x_{10}-4x_{10}x_{57}-x_{19}-4x_{57}^{2}
\\ &+4x_{58}-4x_{58}^{2}+x_{59}-x_{71}
\end{aligned}\]
\[\begin{aligned}
r_{45}={}&-1-x_{1}+x_{6}x_{7}+x_{8}-x_{8}^{2}-x_{8}x_{9}-x_{8}x_{10}+4x_{9}-x_{9}^{2}-x_{9}x_{10}+2x_{10}^{2}
\\ &+4x_{10}x_{57}-x_{19}+4x_{57}^{2}+4x_{58}-4x_{58}^{2}-x_{59}-4x_{59}^{2}+x_{60}-x_{71}
\end{aligned}\]
\[\begin{aligned}
r_{46}={}&-1-x_{1}+x_{6}x_{7}-4x_{60}-4x_{60}^{2}+x_{61}
\end{aligned}\]
\[\begin{aligned}
r_{47}={}&-x_{3}x_{13}+x_{62}
\end{aligned}\]
\[\begin{aligned}
r_{48}={}&x_{12}-x_{13}-2x_{22}x_{62}-x_{63}
\end{aligned}\]
\[\begin{aligned}
r_{49}={}&x_{12}-x_{13}+2x_{22}x_{62}-2x_{26}x_{62}-2x_{30}x_{62}+x_{63}-x_{64}-x_{66}
\end{aligned}\]
\[\begin{aligned}
r_{50}={}&-x_{3}x_{13}-2x_{28}x_{62}+2x_{41}x_{62}-x_{62}-x_{65}+x_{70}
\end{aligned}\]
\[\begin{aligned}
r_{51}={}&-12x_{10}x_{62}+x_{12}-x_{13}-2x_{30}x_{62}-x_{66}
\end{aligned}\]
\[\begin{aligned}
r_{52}={}&x_{12}-x_{13}+2x_{22}x_{62}-2x_{28}x_{62}-2x_{32}x_{62}+x_{63}-x_{65}-x_{67}
\end{aligned}\]
\[\begin{aligned}
r_{53}={}&x_{12}-x_{13}-4x_{34}x_{62}-2x_{35}x_{62}-x_{68}
\end{aligned}\]
\[\begin{aligned}
r_{54}={}&x_{12}-x_{13}-2x_{38}x_{62}-x_{69}
\end{aligned}\]
\[\begin{aligned}
r_{55}={}&-2x_{3}x_{13}+x_{12}-x_{13}-2x_{41}x_{62}-2x_{62}-x_{70}
\end{aligned}\]
\[\begin{aligned}
r_{56}={}&-x_{8}x_{9}+x_{8}x_{10}+x_{9}+x_{9}x_{10}-x_{10}-x_{10}^{2}+x_{71}
\end{aligned}\]
\[\begin{aligned}
r_{57}={}&1+2x_{11}+x_{11}^{2}-x_{71}^{2}+x_{72}
\end{aligned}\]
\[\begin{aligned}
r_{58}={}&x_{3}x_{13}-4x_{60}x_{62}-2x_{61}x_{62}+x_{62}+x_{73}
\end{aligned}\]
\[\begin{aligned}
r_{59}={}&-x_{12}^{2}+2x_{12}x_{13}+2x_{12}x_{73}-x_{13}^{2}-2x_{13}x_{73}-x_{72}^{2}-x_{73}^{2}+x_{74}
\end{aligned}\]
\[\begin{aligned}
r_{60}={}&-1-x_{3}x_{9}+x_{75}
\end{aligned}\]
\[\begin{aligned}
r_{61}={}&-x_{13}x_{75}+x_{77}x_{78}+x_{79}
\end{aligned}\]
\[\begin{aligned}
r_{62}={}&-1+x_{3}-x_{3}x_{8}+x_{76}
\end{aligned}\]
\[\begin{aligned}
r_{63}={}&-1-x_{1}+x_{3}+x_{3}x_{8}+4x_{3}x_{9}-2x_{3}x_{19}+x_{6}x_{7}+2x_{75}-x_{76}+x_{77}
\end{aligned}\]
\[\begin{aligned}
r_{64}={}&-x_{12}x_{76}+x_{78}
\end{aligned}\]
\[\begin{aligned}
r_{65}={}&5x_{12}+x_{12}x_{75}+7x_{13}-5x_{13}x_{75}-4x_{13}x_{76}-x_{13}x_{77}-2x_{22}x_{62}-2x_{26}x_{62}
\\ &-2x_{30}x_{62}-2x_{32}x_{62}-2x_{38}x_{62}-x_{63}-x_{64}-x_{66}-x_{67}-x_{69}-3x_{77}x_{78}-3x_{79}+x_{80}
\end{aligned}\]
\[\begin{aligned}
r_{66}={}&-x_{12}x_{63}+x_{13}x_{63}+x_{81}
\end{aligned}\]
\[\begin{aligned}
r_{67}={}&-x_{64}x_{81}+x_{82}
\end{aligned}\]
\[\begin{aligned}
r_{68}={}&-x_{65}x_{82}+x_{83}
\end{aligned}\]
\[\begin{aligned}
r_{69}={}&-x_{66}x_{83}+x_{84}
\end{aligned}\]
\[\begin{aligned}
r_{70}={}&-x_{67}x_{84}+x_{85}
\end{aligned}\]
\[\begin{aligned}
r_{71}={}&-x_{68}x_{85}+x_{86}
\end{aligned}\]
\[\begin{aligned}
r_{72}={}&-x_{69}x_{86}+x_{87}
\end{aligned}\]
\[\begin{aligned}
r_{73}={}&-x_{70}x_{87}+x_{88}
\end{aligned}\]
\[\begin{aligned}
r_{74}={}&-x_{79}x_{88}+x_{89}
\end{aligned}\]
\[\begin{aligned}
r_{75}={}&-x_{80}x_{89}+x_{90}
\end{aligned}\]
\[\begin{aligned}
r_{76}={}&-x_{90}^{2}+x_{91}
\end{aligned}\]
\[\begin{aligned}
r_{77}={}&-x_{2}-x_{2}x_{16}+x_{16}+x_{92}
\end{aligned}\]
\[\begin{aligned}
r_{78}={}&-x_{14}x_{92}+x_{93}
\end{aligned}\]
\[\begin{aligned}
r_{79}={}&-2x_{1}^{2}+2x_{1}x_{12}+2x_{1}x_{13}+2x_{1}x_{93}+2x_{1}x_{94}-x_{12}^{2}-2x_{12}x_{93}-x_{13}^{2}
\\ &-2x_{13}x_{94}+x_{17}x_{95}-x_{74}x_{91}-x_{93}^{2}-x_{94}^{2}
\end{aligned}\]
\[\begin{aligned}
r_{80}={}&-x_{2}x_{15}+x_{15}-x_{15}x_{92}+x_{94}
\end{aligned}\]
\[\begin{aligned}
r_{81}={}&-x_{98}-4x_{98}x_{230}+x_{100}
\end{aligned}\]
\[\begin{aligned}
r_{82}={}&-x_{4}x_{122}+x_{96}x_{221}-x_{99}-x_{100}x_{122}
\end{aligned}\]
\[\begin{aligned}
r_{83}={}&x_{101}x_{222}-x_{103}x_{125}-x_{104}
\end{aligned}\]
\[\begin{aligned}
r_{84}={}&x_{105}x_{222}-x_{106}x_{125}-x_{107}
\end{aligned}\]
\[\begin{aligned}
r_{85}={}&x_{108}x_{222}-x_{109}x_{125}-x_{110}
\end{aligned}\]
\[\begin{aligned}
r_{86}={}&x_{111}x_{222}-x_{112}x_{125}-x_{113}
\end{aligned}\]
\[\begin{aligned}
r_{87}={}&x_{114}x_{223}-x_{116}x_{125}-x_{117}
\end{aligned}\]
\[\begin{aligned}
r_{88}={}&x_{118}x_{223}-x_{119}x_{125}-x_{120}
\end{aligned}\]
\[\begin{aligned}
r_{89}={}&-x_{17}x_{122}+x_{125}
\end{aligned}\]
\[\begin{aligned}
r_{90}={}&x_{121}-x_{123}x_{125}-x_{124}
\end{aligned}\]
\[\begin{aligned}
r_{91}={}&-x_{125}x_{127}+x_{126}-x_{128}
\end{aligned}\]
\[\begin{aligned}
r_{92}={}&-x_{125}x_{130}+x_{129}x_{221}-x_{131}
\end{aligned}\]
\[\begin{aligned}
r_{93}={}&-x_{125}x_{133}+x_{132}x_{222}-x_{134}
\end{aligned}\]
\[\begin{aligned}
r_{94}={}&x_{135}-2x_{135}x_{137}+x_{138}
\end{aligned}\]
\[\begin{aligned}
r_{95}={}&-x_{135}x_{138}-x_{138}+x_{139}
\end{aligned}\]
\[\begin{aligned}
r_{96}={}&x_{138}^{2}-x_{139}^{2}+x_{159}
\end{aligned}\]
\[\begin{aligned}
r_{97}={}&-x_{136}x_{141}-x_{141}+x_{142}
\end{aligned}\]
\[\begin{aligned}
r_{98}={}&2x_{140}+x_{140}^{2}+x_{141}^{2}-x_{142}^{2}
\end{aligned}\]
\[\begin{aligned}
r_{99}={}&1-2x_{0}+x_{0}^{2}-2x_{0}x_{136}+x_{143}
\end{aligned}\]
\[\begin{aligned}
r_{100}={}&-1-x_{0}x_{141}+x_{3}-x_{140}+x_{142}+x_{143}x_{144}
\end{aligned}\]
\[\begin{aligned}
r_{101}={}&1-2x_{4}+x_{4}^{2}-2x_{4}x_{136}+x_{145}
\end{aligned}\]
\[\begin{aligned}
r_{102}={}&-x_{4}x_{190}+x_{17}+x_{145}x_{146}-x_{189}+x_{191}
\end{aligned}\]
\[\begin{aligned}
r_{103}={}&-1+2x_{136}x_{205}-x_{147}+2x_{205}-x_{205}^{2}
\end{aligned}\]
\[\begin{aligned}
r_{104}={}&1+x_{147}x_{148}-x_{189}-x_{190}x_{205}+x_{191}+4x_{230}
\end{aligned}\]
\[\begin{aligned}
r_{105}={}&-x_{136}x_{150}-x_{150}+x_{151}
\end{aligned}\]
\[\begin{aligned}
r_{106}={}&2x_{149}+x_{149}^{2}+x_{150}^{2}-x_{151}^{2}
\end{aligned}\]
\[\begin{aligned}
r_{107}={}&-x_{136}+x_{136}^{2}-8x_{136}x_{230}-2x_{147}x_{148}+x_{152}-x_{159}-x_{189}^{2}-x_{190}^{2}
\\ &+2x_{190}x_{205}+x_{191}^{2}+x_{207}x_{214}-x_{211}+16x_{230}^{2}
\end{aligned}\]
\[\begin{aligned}
r_{108}={}&-1+x_{96}+x_{136}x_{150}-x_{149}-x_{150}-4x_{150}x_{230}+x_{152}x_{153}
\end{aligned}\]
\[\begin{aligned}
r_{109}={}&-1+x_{97}+x_{136}x_{141}-2x_{136}x_{154}-x_{140}-4x_{141}x_{230}+x_{152}x_{154}+x_{154}
\\ &+8x_{154}x_{230}
\end{aligned}\]
\[\begin{aligned}
r_{110}={}&1+x_{136}-x_{136}x_{156}-x_{156}+x_{157}
\end{aligned}\]
\[\begin{aligned}
r_{111}={}&-x_{136}+x_{136}x_{156}+2x_{155}+x_{155}^{2}-x_{156}+x_{156}^{2}-x_{157}-x_{157}^{2}
\end{aligned}\]
\[\begin{aligned}
r_{112}={}&-1+x_{101}-x_{136}+x_{136}x_{156}-x_{147}x_{148}+x_{152}x_{158}-x_{155}-x_{156}
\\ &-4x_{156}x_{230}+x_{189}+x_{190}x_{205}-x_{191}
\end{aligned}\]
\[\begin{aligned}
r_{113}={}&-x_{136}-x_{136}^{2}+x_{159}+x_{207}x_{214}-x_{211}
\end{aligned}\]
\[\begin{aligned}
r_{114}={}&x_{102}+x_{136}x_{142}-2x_{136}x_{160}-x_{140}-4x_{140}x_{230}-x_{141}x_{159}-4x_{142}x_{230}
\\ &+x_{147}x_{148}+x_{152}x_{160}+x_{160}+8x_{160}x_{230}-x_{189}-x_{190}x_{205}+x_{191}
\end{aligned}\]
\[\begin{aligned}
r_{115}={}&-x_{136}x_{162}-x_{162}+x_{163}
\end{aligned}\]
\[\begin{aligned}
r_{116}={}&2x_{161}+x_{161}^{2}+x_{162}^{2}-x_{163}^{2}
\end{aligned}\]
\[\begin{aligned}
r_{117}={}&-1+x_{105}+x_{136}x_{162}+x_{152}x_{164}-x_{161}-x_{162}-4x_{162}x_{230}
\end{aligned}\]
\[\begin{aligned}
r_{118}={}&-x_{136}x_{166}-x_{166}+x_{167}
\end{aligned}\]
\[\begin{aligned}
r_{119}={}&2x_{165}+x_{165}^{2}+x_{166}^{2}-x_{167}^{2}
\end{aligned}\]
\[\begin{aligned}
r_{120}={}&-1+x_{108}+x_{136}x_{166}+x_{152}x_{168}-x_{165}-x_{166}-4x_{166}x_{230}
\end{aligned}\]
\[\begin{aligned}
r_{121}={}&-x_{136}x_{170}-x_{170}+x_{171}
\end{aligned}\]
\[\begin{aligned}
r_{122}={}&2x_{169}+x_{169}^{2}+x_{170}^{2}-x_{171}^{2}
\end{aligned}\]
\[\begin{aligned}
r_{123}={}&-1+x_{111}+x_{136}x_{170}+x_{152}x_{172}-x_{169}-x_{170}-4x_{170}x_{230}
\end{aligned}\]
\[\begin{aligned}
r_{124}={}&-x_{136}x_{174}-x_{174}+x_{175}
\end{aligned}\]
\[\begin{aligned}
r_{125}={}&2x_{173}+x_{173}^{2}+x_{174}^{2}-x_{175}^{2}
\end{aligned}\]
\[\begin{aligned}
r_{126}={}&-1+x_{114}+x_{136}x_{174}+x_{152}x_{176}-x_{173}-x_{174}-4x_{174}x_{230}
\end{aligned}\]
\[\begin{aligned}
r_{127}={}&-x_{136}x_{178}-x_{178}+x_{179}
\end{aligned}\]
\[\begin{aligned}
r_{128}={}&2x_{177}+x_{177}^{2}+x_{178}^{2}-x_{179}^{2}
\end{aligned}\]
\[\begin{aligned}
r_{129}={}&-1+x_{115}+x_{136}x_{178}-2x_{136}x_{180}+x_{152}x_{180}-x_{177}-4x_{178}x_{230}+x_{180}
\\ &+8x_{180}x_{230}
\end{aligned}\]
\[\begin{aligned}
r_{130}={}&-x_{136}x_{182}-x_{182}+x_{183}
\end{aligned}\]
\[\begin{aligned}
r_{131}={}&2x_{181}+x_{181}^{2}+x_{182}^{2}-x_{183}^{2}
\end{aligned}\]
\[\begin{aligned}
r_{132}={}&-1+x_{118}+x_{136}x_{182}+x_{152}x_{184}-x_{181}-x_{182}-4x_{182}x_{230}
\end{aligned}\]
\[\begin{aligned}
r_{133}={}&-x_{136}x_{186}-x_{186}+x_{187}
\end{aligned}\]
\[\begin{aligned}
r_{134}={}&2x_{185}+x_{185}^{2}+x_{186}^{2}-x_{187}^{2}
\end{aligned}\]
\[\begin{aligned}
r_{135}={}&-1+x_{121}+x_{136}x_{186}+x_{152}x_{188}-x_{185}-x_{186}-4x_{186}x_{230}
\end{aligned}\]
\[\begin{aligned}
r_{136}={}&1-x_{136}x_{190}-x_{190}+x_{191}
\end{aligned}\]
\[\begin{aligned}
r_{137}={}&-1-x_{4}x_{190}+x_{17}+x_{145}x_{146}+x_{189}+x_{189}^{2}+x_{190}^{2}-x_{191}-x_{191}^{2}
\end{aligned}\]
\[\begin{aligned}
r_{138}={}&x_{122}-2x_{136}x_{192}+x_{152}x_{192}-x_{189}-x_{190}-4x_{190}x_{230}+x_{191}+x_{192}
\\ &+8x_{192}x_{230}
\end{aligned}\]
\[\begin{aligned}
r_{139}={}&-x_{136}x_{194}-x_{194}+x_{195}
\end{aligned}\]
\[\begin{aligned}
r_{140}={}&2x_{193}+x_{193}^{2}+x_{194}^{2}-x_{195}^{2}
\end{aligned}\]
\[\begin{aligned}
r_{141}={}&-1+x_{126}+x_{136}x_{194}+x_{152}x_{196}-x_{193}-x_{194}-4x_{194}x_{230}
\end{aligned}\]
\[\begin{aligned}
r_{142}={}&-x_{136}x_{198}-x_{198}+x_{199}
\end{aligned}\]
\[\begin{aligned}
r_{143}={}&2x_{197}+x_{197}^{2}+x_{198}^{2}-x_{199}^{2}
\end{aligned}\]
\[\begin{aligned}
r_{144}={}&-1+x_{129}+x_{136}x_{198}+x_{152}x_{200}-x_{197}-x_{198}-4x_{198}x_{230}
\end{aligned}\]
\[\begin{aligned}
r_{145}={}&-x_{136}x_{202}-x_{202}+x_{203}
\end{aligned}\]
\[\begin{aligned}
r_{146}={}&2x_{201}+x_{201}^{2}+x_{202}^{2}-x_{203}^{2}
\end{aligned}\]
\[\begin{aligned}
r_{147}={}&-1+x_{132}+x_{136}x_{202}+x_{152}x_{204}-x_{201}-x_{202}-4x_{202}x_{230}
\end{aligned}\]
\[\begin{aligned}
r_{148}={}&-1-2x_{0}-4x_{0}^{2}-2x_{0}x_{3}-x_{1}-x_{1}^{2}-x_{2}^{2}-2x_{3}^{2}-x_{4}^{2}+x_{4}x_{190}+x_{6}x_{7}+x_{8}
\\ &+2x_{8}x_{9}+x_{9}-x_{10}^{2}-x_{11}^{2}-x_{12}^{2}-x_{13}^{2}-x_{14}^{2}-x_{15}^{2}-x_{16}^{2}-x_{17}-x_{17}^{2}-x_{19}-x_{96}^{2}
\\ &-x_{97}^{2}-x_{101}^{2}-x_{102}^{2}-x_{105}^{2}-x_{108}^{2}-x_{111}^{2}-x_{114}^{2}-x_{115}^{2}-x_{118}^{2}-x_{121}^{2}-x_{122}^{2}
\\ &-x_{126}^{2}-x_{129}^{2}-x_{132}^{2}+x_{135}-2x_{136}+2x_{136}^{2}-x_{136}x_{156}-16x_{136}x_{230}+x_{138}^{2}
\\ &-x_{139}^{2}+x_{141}^{2}-x_{145}x_{146}+2x_{147}x_{148}+x_{150}^{2}+2x_{152}-2x_{155}-x_{155}^{2}-x_{156}
\\ &+x_{157}+x_{157}^{2}-x_{159}+x_{162}^{2}+x_{166}^{2}+x_{170}^{2}+x_{174}^{2}+x_{178}^{2}+x_{182}^{2}+x_{186}^{2}-x_{189}
\\ &+x_{190}^{2}-2x_{190}x_{205}+x_{191}+x_{194}^{2}+x_{198}^{2}-x_{202}^{2}-x_{205}^{2}+x_{207}x_{214}-x_{211}
\\ &+x_{224}^{2}+x_{225}^{2}+x_{226}+x_{226}^{2}
\end{aligned}\]
\[\begin{aligned}
r_{149}={}&-x_{4}-x_{4}x_{8}-x_{4}x_{9}-x_{4}x_{10}-x_{4}x_{11}-x_{4}x_{12}-x_{4}x_{13}-x_{4}x_{14}-x_{4}x_{15}
\\ &-x_{4}x_{16}+x_{205}
\end{aligned}\]
\[\begin{aligned}
r_{150}={}&-2+x_{0}x_{2}+x_{0}x_{141}-2x_{1}-x_{2}+x_{2}x_{5}+x_{3}x_{8}+x_{3}x_{9}+2x_{3}x_{231}-x_{4}x_{18}
\\ &+x_{6}x_{7}-x_{7}+x_{9}^{2}+2x_{9}x_{10}+x_{10}+x_{10}^{2}+x_{16}-x_{23}-x_{75}-x_{76}-x_{136}x_{190}+x_{140}
\\ &-x_{142}-x_{143}x_{144}-x_{190}+x_{191}+x_{201}-2x_{206}^{2}-x_{207}^{2}+x_{207}x_{215}-x_{208}^{2}-x_{209}
\\ &-x_{209}^{2}-x_{210}^{2}+x_{217}^{2}+x_{218}^{2}
\end{aligned}\]
\[\begin{aligned}
r_{151}={}&-x_{136}x_{208}-x_{208}+x_{217}
\end{aligned}\]
\[\begin{aligned}
r_{152}={}&-1+x_{207}^{2}+x_{208}^{2}-x_{217}^{2}
\end{aligned}\]
\[\begin{aligned}
r_{153}={}&-x_{210}-x_{210}x_{211}+x_{218}
\end{aligned}\]
\[\begin{aligned}
r_{154}={}&-1+x_{209}^{2}+x_{210}^{2}-x_{218}^{2}
\end{aligned}\]
\[\begin{aligned}
r_{155}={}&-x_{202}^{2}+x_{219}
\end{aligned}\]
\[\begin{aligned}
r_{156}={}&x_{208}-2x_{213}x_{219}-x_{219}
\end{aligned}\]
\[\begin{aligned}
r_{157}={}&x_{136}x_{208}-3x_{208}-8x_{208}x_{212}+x_{211}-x_{217}
\end{aligned}\]
\[\begin{aligned}
r_{158}={}&-x_{136}-x_{207}x_{214}+x_{211}
\end{aligned}\]
\[\begin{aligned}
r_{159}={}&-1-x_{201}-x_{207}x_{215}+x_{209}
\end{aligned}\]
\[\begin{aligned}
r_{160}={}&-1-4x_{8}-x_{8}^{2}-x_{8}x_{9}-x_{8}x_{10}+x_{9}-x_{9}^{2}-x_{9}x_{10}+x_{10}+x_{10}^{2}-3x_{19}-2x_{22}
\\ &-2x_{26}-2x_{28}-2x_{30}-2x_{32}-4x_{34}-2x_{35}-2x_{38}-2x_{41}-4x_{60}-2x_{61}-x_{71}
\\ &-4x_{202}x_{216}+x_{210}
\end{aligned}\]
\[\begin{aligned}
r_{161}={}&-1+2x_{0}^{2}+2x_{0}x_{3}-2x_{0}x_{141}-2x_{0}x_{150}+x_{1}^{2}-2x_{1}x_{190}+x_{2}^{2}-2x_{2}x_{178}
\\ &+x_{3}^{2}-2x_{3}x_{150}+x_{8}^{2}-2x_{8}x_{156}-10x_{8}x_{202}+x_{9}^{2}-2x_{9}x_{162}+6x_{9}x_{202}+x_{10}^{2}
\\ &-2x_{10}x_{166}+x_{11}^{2}-2x_{11}x_{170}+x_{12}^{2}-2x_{12}x_{174}+x_{13}^{2}-2x_{13}x_{182}+x_{14}^{2}
\\ &-2x_{14}x_{186}+x_{15}^{2}-2x_{15}x_{194}+x_{16}^{2}-2x_{16}x_{198}-8x_{19}x_{202}-4x_{22}x_{202}
\\ &-4x_{26}x_{202}-4x_{28}x_{202}-4x_{30}x_{202}-4x_{32}x_{202}-8x_{34}x_{202}-4x_{35}x_{202}
\\ &-4x_{38}x_{202}-4x_{41}x_{202}-8x_{60}x_{202}-4x_{61}x_{202}+x_{136}-x_{136}x_{156}-x_{136}x_{220}
\\ &+x_{140}^{2}+x_{149}^{2}+x_{155}^{2}+x_{156}+x_{157}+x_{161}^{2}+x_{165}^{2}+x_{169}^{2}+x_{173}^{2}+x_{177}^{2}+x_{181}^{2}
\\ &+x_{185}^{2}+x_{189}^{2}+x_{193}^{2}+x_{197}^{2}+x_{201}^{2}+x_{202}^{2}+x_{219}-x_{224}^{2}-x_{225}^{2}-x_{226}-x_{226}^{2}
\end{aligned}\]
\[\begin{aligned}
r_{162}={}&-x_{97}x_{221}+x_{122}
\end{aligned}\]
\[\begin{aligned}
r_{163}={}&-x_{102}x_{222}+x_{122}
\end{aligned}\]
\[\begin{aligned}
r_{164}={}&-x_{115}x_{223}+x_{122}
\end{aligned}\]
\[\begin{aligned}
r_{165}={}&-1+x_{2}-x_{3}x_{231}
\end{aligned}\]
\[\begin{aligned}
r_{166}={}&1+x_{1}+x_{2}-x_{3}x_{9}-x_{3}x_{231}-x_{6}x_{7}+2x_{8}-6x_{8}^{2}+2x_{8}x_{9}+2x_{8}x_{10}-8x_{8}x_{19}
\\ &-20x_{8}x_{22}-20x_{8}x_{26}-20x_{8}x_{28}-20x_{8}x_{30}-20x_{8}x_{32}-40x_{8}x_{34}-20x_{8}x_{35}
\\ &-20x_{8}x_{38}-20x_{8}x_{41}-40x_{8}x_{60}-20x_{8}x_{61}+2x_{9}+6x_{9}^{2}-4x_{9}x_{10}-8x_{9}x_{19}
\\ &+12x_{9}x_{22}+12x_{9}x_{26}+12x_{9}x_{28}+12x_{9}x_{30}+12x_{9}x_{32}+20x_{9}x_{34}+12x_{9}x_{35}
\\ &+12x_{9}x_{38}+12x_{9}x_{41}+24x_{9}x_{60}+12x_{9}x_{61}-2x_{10}-x_{10}^{2}+8x_{19}-16x_{19}x_{22}
\\ &-16x_{19}x_{26}-16x_{19}x_{28}-16x_{19}x_{30}-16x_{19}x_{32}-32x_{19}x_{34}-16x_{19}x_{35}
\\ &-16x_{19}x_{38}-16x_{19}x_{41}-32x_{19}x_{60}-16x_{19}x_{61}-2x_{22}-4x_{22}^{2}-8x_{22}x_{26}
\\ &-8x_{22}x_{28}-8x_{22}x_{30}-8x_{22}x_{32}-16x_{22}x_{34}-8x_{22}x_{35}-8x_{22}x_{38}-8x_{22}x_{41}
\\ &-16x_{22}x_{60}-8x_{22}x_{61}+3x_{23}-3x_{24}-2x_{26}-4x_{26}^{2}-8x_{26}x_{28}-8x_{26}x_{30}-8x_{26}x_{32}
\\ &-16x_{26}x_{34}-8x_{26}x_{35}-8x_{26}x_{38}-8x_{26}x_{41}-16x_{26}x_{60}-8x_{26}x_{61}-2x_{28}-4x_{28}^{2}
\\ &-8x_{28}x_{30}-8x_{28}x_{32}-16x_{28}x_{34}-8x_{28}x_{35}-8x_{28}x_{38}-8x_{28}x_{41}-16x_{28}x_{60}
\\ &-8x_{28}x_{61}-2x_{30}-4x_{30}^{2}-8x_{30}x_{32}-16x_{30}x_{34}-8x_{30}x_{35}-8x_{30}x_{38}-8x_{30}x_{41}
\\ &-16x_{30}x_{60}-8x_{30}x_{61}-4x_{31}-4x_{31}^{2}-x_{32}-4x_{32}^{2}-16x_{32}x_{34}-8x_{32}x_{35}
\\ &-8x_{32}x_{38}-8x_{32}x_{41}-16x_{32}x_{60}-8x_{32}x_{61}-4x_{34}-20x_{34}^{2}-16x_{34}x_{35}
\\ &-16x_{34}x_{38}-16x_{34}x_{41}-32x_{34}x_{60}-16x_{34}x_{61}-x_{35}-4x_{35}^{2}-8x_{35}x_{38}-8x_{35}x_{41}
\\ &-16x_{35}x_{60}-8x_{35}x_{61}-x_{36}-2x_{38}-4x_{38}^{2}-8x_{38}x_{41}-16x_{38}x_{60}-8x_{38}x_{61}-2x_{41}
\\ &-4x_{41}^{2}-16x_{41}x_{60}-8x_{41}x_{61}+4x_{60}-8x_{60}^{2}-16x_{60}x_{61}-4x_{61}-4x_{61}^{2}+x_{75}
\\ &-x_{141}^{2}-x_{150}^{2}+2x_{156}-x_{156}^{2}-x_{162}^{2}-x_{166}^{2}-x_{170}^{2}-x_{174}^{2}-x_{178}^{2}-x_{182}^{2}
\\ &-x_{186}^{2}-x_{190}^{2}-x_{194}^{2}-x_{198}^{2}+x_{202}^{2}-4x_{202}x_{216}+x_{210}-x_{224}^{2}-x_{225}^{2}-x_{226}
\\ &-x_{226}^{2}
\end{aligned}\]
\[\begin{aligned}
r_{167}={}&-2x_{99}^{2}-2x_{104}^{2}-2x_{107}^{2}-2x_{110}^{2}-2x_{113}^{2}-2x_{117}^{2}-2x_{120}^{2}+x_{122}^{2}
\\ &-2x_{124}^{2}-2x_{128}^{2}-2x_{131}^{2}-2x_{134}^{2}-x_{227}^{2}-x_{228}^{2}-x_{229}-x_{229}^{2}
\end{aligned}\]
\[\begin{aligned}
r_{168}={}&-1+x_{231}-4x_{232}-4x_{232}^{2}
\end{aligned}\]
\endgroup


\begin{thebibliography}{99}
\bibitem{current} V. Reshetnikov et al., \emph{ProveIt: DiophantineConsistency},
commit \path{fb1bed5eb7196cf35abfaa746c18c6ee7eae6cfb},
\texttt{CURRENT\_EQUATION.json} and \texttt{equation.tex}.
\url{https://github.com/VladimirReshetnikov/ProveIt/tree/fb1bed5eb7196cf35abfaa746c18c6ee7eae6cfb/Computability/DiophantineConsistency}.
\bibitem{native} ProveIt, \emph{A zero-weight physical flag and shorter separated clause blocks},
\texttt{research/zero-weight-flag-independent/PROOF.md}, same commit.
The frozen text and numerical source data are included in this release.
\bibitem{trace} ProveIt package 163, \emph{Use the explicit Pell trace before squaring it},
\texttt{163-trace-radical/compact/PROOF.md}, same commit.
\bibitem{inverse} ProveIt package 164, \emph{Inverse-Lucas correction with projection before denominator recovery},
\texttt{164-unsquared-factors-linear-selectors/positive/PROOF.md};
and its signed linear-selector proof, same commit.
\bibitem{wide} ProveIt package 165, \emph{Wider budgets, round 20},
\texttt{165-wider-budgets-round20/report/power\_base.tex},
\texttt{report/signed\_new.tex}, and the literal
\texttt{signed/power239/equation.json}, same commit.
\bibitem{radical} ProveIt package 159, \emph{A direct quartic for Con(ZFC) in 234 integer variables},
\texttt{159-quadratic-projection-234}, same commit.
\bibitem{threesquares} A. Danilkin and L. Chevalier,
\emph{Three Squares Theorem}, Archive of Formal Proofs, 3 May 2023.
\url{https://isa-afp.org/entries/Three_Squares.html}.
\end{thebibliography}
\end{document}
